\documentclass[11pt]{article}
\usepackage[T1]{fontenc}
\usepackage[utf8]{inputenc}
\usepackage{lmodern,microtype}
\usepackage[margin=1in,headheight=15pt]{geometry}
\usepackage{amsmath,amssymb,amsthm,mathtools}
\usepackage{booktabs,longtable,array,enumitem}
\usepackage{xcolor,fancyhdr}
\usepackage[hidelinks]{hyperref}
\hypersetup{pdftitle={The Gamma Function in Intuitionistic Primitive Recursive Arithmetic},pdfauthor={Prepared for Vladimir Reshetnikov},pdfsubject={Constructive finite-precision proofs and a Lean formalization roadmap}}
\usepackage[nameinlink,noabbrev]{cleveref}
\definecolor{ink}{RGB}{24,48,70}
\pagestyle{fancy}\fancyhf{}\fancyhead[L]{\small Gamma in intuitionistic primitive recursive arithmetic}\fancyhead[R]{\small\thepage}
\setlength{\parskip}{0.35em}
\setlength{\emergencystretch}{2em}
\newtheorem{theorem}{Theorem}[section]
\newtheorem{lemma}[theorem]{Lemma}
\newtheorem{proposition}[theorem]{Proposition}
\newtheorem{corollary}[theorem]{Corollary}
\theoremstyle{definition}
\newtheorem{definition}[theorem]{Definition}
\newtheorem{construction}[theorem]{Construction}
\newtheorem{remark}[theorem]{Remark}
\newcommand{\N}{\mathbb N}\newcommand{\Z}{\mathbb Z}\newcommand{\Q}{\mathbb Q}\newcommand{\R}{\mathbb R}\newcommand{\C}{\mathbb C}
\newcommand{\ipra}{\ensuremath{\mathsf{iPRA}}}\newcommand{\PR}{\ensuremath{\mathsf{PR}}}
\newcommand{\abs}[1]{\left|#1\right|}\newcommand{\norm}[1]{\left\|#1\right\|}
\newcommand{\ceil}[1]{\left\lceil#1\right\rceil}
\newcommand{\dd}{\,\mathrm d}\newcommand{\e}{\mathrm e}
\newcommand{\Rea}{\operatorname{Re}}\newcommand{\Ima}{\operatorname{Im}}
\newcommand{\Res}{\operatorname{Res}}\newcommand{\sgn}{\operatorname{sgn}}
\newcommand{\code}[1]{\texttt{#1}}
\newcommand{\Rising}[2]{(#1)_{#2}}
\title{\color{ink}The Gamma Function in\, Intuitionistic\, Primitive Recursive Arithmetic\\[0.5em]\large Finite products, effective integrals, functional identities,\newline derivatives, convexity, and asymptotic estimates}
\author{A research development prepared for Vladimir Reshetnikov}
\date{15 November 2026}
\begin{document}
\maketitle
\begin{abstract}
We give a constructive arithmetic development of a substantial collection of
gamma-function statements appearing in the two requested reference articles.
The main objects are fixed primitive recursive rational approximations with
explicit error certificates. Euler's finite products, rather than an
unrestricted real-completeness axiom, construct the function. We prove its
integral representation, functional equation, factorial and half-integer
values, entire reciprocal and simple poles, reflection and multiplication
formulas, beta identities, differential and local power-series formulas,
log-convexity and effective Bohr--Mollerup uniqueness, real Stirling estimates,
and several further integral and special-value consequences. Each infinite
notation has a finite-precision interpretation. The report distinguishes
these human-readable arithmetic proofs from the accompanying Lean work:
the executable definitions and elementary finite lemmas are checked, while
the full deep embedding of the analytic estimates remains a specified
formalization project. No completed Lean proof of the gamma theory is claimed.
\end{abstract}
\tableofcontents
\newpage
\section{The precise constructive claim}
\label{sec:foundations}

The setting is first-order intuitionistic primitive recursive arithmetic,
with a symbol for each primitive recursive definition, its defining
equations, successor separation and injectivity, intuitionistic predicate
logic with equality, and induction for quantifier-free formulas. This is
the system \code{PRAAx} in the inspected \code{ConciseIPRA/PR.lean}.
Its arithmetic interpretation in the native concise calculus is supplied by
\code{pra\_iff\_native} in \code{NativeBiInterpretation.lean}~\cite{repo}.
We work in this ordinary presentation to make the mathematics readable.

There is no sort of arbitrary real numbers or arbitrary real functions in
this theory. Accordingly, an assertion about a real or complex argument
means a theorem schema for \emph{each fixed primitive recursive name}, with
the displayed numerical domain information and error certificates. The
rational-argument estimates are uniform first-order statements about
rational codes and precisions. Substituting a fixed name and its modulus
gives the corresponding theorem about that name. We do not postulate a
primitive recursive universal evaluator for all primitive recursive
programs; no such evaluator exists.

\subsection{Finite arithmetic and names}

Integers are sign-and-magnitude codes, and a rational is a pair $(u,v)$
with $u\in\Z$ and $v\in\N_{>0}$. Equality and order are cross-multiplication
tests. Reduction by gcd is convenient but unnecessary. Finite lists,
factorials, powers, sums, products, finite differences, and bounded search
are primitive recursive. Complex rationals are pairs of rationals. For
coding we use $\norm{x+iy}_1=|x|+|y|$; estimates written with the ordinary
complex modulus can instead be checked after squaring, or transferred with
$|z|\le\norm z_1\le2|z|$. Enlarge a disk radius by a factor of two when
using the latter convention.

\begin{definition}[Names and equality]
A regular real name is a fixed primitive recursive map $a:\N\to\Q$
satisfying
\[
 |a(p)-a(q)|\le2^{-p}+2^{-q}\qquad(p,q\in\N).
\]
A complex name is a pair of such names; alternatively use this inequality
in the $1$-norm. Equality of names is
\[
 a\simeq b\quad\Longleftrightarrow\quad
 \forall p\; |a(p+2)-b(p+2)|\le2^{-p}.
\]
An inequality $a\le b$ means the analogous comparison with error
$2^{-p}$. Positive separation from zero is supplied as a rational number
$d>0$ and the certificate $a\ge d$. Complex apartness is supplied as
$|a|\ge d$; it is not inferred from $\neg(a=0)$.
\end{definition}

These equalities are equivalence relations. For example, to prove
transitivity at precision $p$, use the hypotheses at precision $p+r$,
regularity to compare stages, and take $r\ge4$; to remove a residual
positive rational error, assume the desired rational inequality fails and
choose a precision bounded by the binary length of its positive gap.
This is a bounded rational argument, not excluded middle for equality of
names. Equivalently, one may define equality by
$\forall p,q\;|a(p)-b(q)|\le2^{-p}+2^{-q}$; the same refinement argument
proves the equivalence of the conventions.

A uniformly convergent finite expression $F_N$ is represented by
\emph{substitution of a modulus}: if
$|F_M-F_N|\le C/N$ for $M\ge N\ge N_0$, use
$N(p)=N_0+\ceil C\,2^{p+4}$ and evaluate $F_{N(p)}$ to error
$2^{-p-4}$. The output is a regular name after a fixed precision shift.
This defines an actual primitive recursive function. It does not invoke
a principle saying that an arbitrary Cauchy sequence has a limit.

\subsection{How the displayed proofs become arithmetic proofs}

Every finite sum identity is proved by induction on its finite length.
Every denominator has a displayed positive lower bound. A bounded
quantifier over a decidable predicate is represented by a finite Boolean
product or sum, hence by a primitive recursive characteristic function.
Thus the invariant to which induction is applied can be an equation
between primitive recursive terms. When several variables or derivative
orders are present, they remain parameters of this quantifier-free
induction.

There is one essential discipline. We never use induction on a statement
such as ``two names are equal,'' which already contains an unbounded
precision quantifier. We first fix the requested precision, replace every
name in the finite calculation by a rational approximation, and allocate
its error budget using the following elementary estimates:
\begin{align}
 \left|\sum_{j<s}a_j-\sum_{j<s}b_j\right|&\le s\varepsilon,
 \label{eq:sum-budget}\\
 \left|\prod_{j<s}a_j-\prod_{j<s}b_j\right|
 &\le sM^{s-1}\varepsilon,
 \quad |a_j|,|b_j|\le M,
 \label{eq:product-budget}\\
 |a^{-1}-b^{-1}|&\le\varepsilon/d^2,
 \quad |a|,|b|\ge d.\label{eq:inverse-budget}
\end{align}
Here $|a_j-b_j|\le\varepsilon$. The product estimate follows by replacing
one factor at a time, so its proof is a finite telescoping identity. When
only $|a|\ge d$ is known, evaluate $b$ with error at most $d/2$ and use
$2\varepsilon/d^2$ instead. Integer upper bounds for all constants suffice.
For an integer $C\ge1$, the unnecessarily large shift $C+k+3$ makes
$C2^{-(C+k+3)}<2^{-k}$, eliminating any need for logarithmic search.

\begin{proposition}[Finite-precision reading]
The proofs below give, for each fixed tuple of names and each requested
precision, finite rational identities or inequalities with primitive
recursive lengths, cutoffs, and precision shifts. Their induction
invariants are quantifier-free. Consequently their displayed equality,
inequality, derivative, and integral conclusions have derivations in
$\ipra$ under the stated certificates.
\end{proposition}

\begin{proof}
Expand rational operations into integer operations, replace finite sums
and products by their recursors, and use
\eqref{eq:sum-budget}--\eqref{eq:inverse-budget} after each arithmetic
operation. All searches in the constructions have an explicit finite
length. The only case distinctions are on integers or rational
comparisons. A series tail is always replaced by a displayed bound before
the target precision is chosen. An integral is always replaced by a
finite sum before invoking a calculus identity. The specific invariants
and moduli are collected in \cref{sec:formalization}; the estimates are
proved in their mathematical sections. Finally discharge the precision
parameter by intuitionistic universal introduction. This is a proof
translation recipe for these concrete arguments, not a claim that
primitive recursiveness alone implies provability of arbitrary statements.
\end{proof}

\subsection{Elementary functions with explicit algorithms}

For $|w|\le B$, $B\ge1$ an integer, define $\exp w$ by the finite Taylor
polynomials. If $m+1\ge2B$, the remaining terms have successive absolute
ratios at most $1/2$, and therefore
\begin{equation}
 \left|\exp w-\sum_{j=0}^{m}\frac{w^j}{j!}\right|
 \le\frac{2B^{m+1}}{(m+1)!}.\label{eq:exp-tail}
\end{equation}
For example $m=4B+p+6$ makes this less than $2^{-p}$: the term at
$2B$ is at most $3^B$, since finite multinomial expansion and
$\sum_{j\ge0}1/j!<3$ give $\sum B^j/j!\le3^B$, and the subsequent
ratios are at most $1/2$. The bound $\sum1/j!<3$ itself follows from
$j!\ge2^{j-1}$ for $j\ge1$. All these are finite inequalities with
geometric remainder bounds. Finite binomial convolution proves
$\exp(u+v)=\exp u\exp v$. In particular $\exp u\ne0$ with explicit
lower bound $3^{-B}$ for $|u|\le B$.

For $q>0$, put $v=(q-1)/(q+1)$. Define
\begin{equation}
 \log q=2\sum_{j\ge0}\frac{v^{2j+1}}{2j+1},\qquad
 |\text{tail after }j<L|\le
 \frac{2|v|^{2L+1}}{(2L+1)(1-v^2)}.\label{eq:log-alg}
\end{equation}
If $B^{-1}\le q\le B$, then $|v|\le1-1/B$ and
$(1-1/B)^B\le1/2$ (use the finite binomial inequality on the
reciprocal). Thus $L=B(p+2B+4)$ is a safe primitive recursive cutoff.
Finite differentiation of these series gives $(\log q)'=1/q$,
$(\exp q)'=\exp q$; the derivative tails are geometric on slightly
smaller intervals. The product and inverse identities, and
$\log(ab)=\log a+\log b$ for positive $a,b$, follow either by finite
series convolution or by integrating their zero derivative on a compact
interval. Rational or complex powers of a \emph{positive real} $t$ mean
$t^z=\exp(z\log t)$. This convention has no complex branch ambiguity.

Define sine and cosine by their Taylor series, equivalently by
$(\exp(iw)-\exp(-iw))/(2i)$ and
$(\exp(iw)+\exp(-iw))/2$. The addition identities follow from the
exponential identity. Define
$\pi=4\int_0^1(1+t^2)^{-1}\dd t$; the alternating arctangent sum has
error at most $4/(2L+1)$, so $L=2^{p+4}$ suffices.
For completeness the identification with trigonometric period is finite:
write $a(u)=\int_0^u(1+t^2)^{-1}\dd t$, differentiate and check
\[
 \sin(2a(u))=\frac{2u}{1+u^2},\qquad
 \cos(2a(u))=\frac{1-u^2}{1+u^2}.
\]
Both pairs have the same initial values and satisfy the same two first
order equations. Uniqueness here requires no general differential
equation theorem: the squared difference has zero derivative, and the
finite fundamental theorem described below makes it zero. At $u=1$ this
gives $\sin(\pi/2)=1$, $\cos(\pi/2)=0$ and then
$\sin\pi=0$, $\cos\pi=-1$. The map $u\mapsto2a(u)$ has derivative
at least one on $[0,1]$. The approximate-sign rational bisection described
in the trigonometric section therefore constructs its inverse on
$[0,\pi/2]$ with a prescribed rate. Positivity on the first quadrant
follows from this parametrization. Concavity, $\sin''=-\sin\le0$, gives
\begin{equation}
 \sin v\ge2v/\pi\quad(0\le v\le\pi/2).\label{eq:sine-chord}
\end{equation}
These facts, addition and periodicity give the needed finite lists of
trigonometric roots. The coarse rational bounds $2<\pi<4$ follow directly
from the defining integral.

\subsection{The finite calculus convention}
\label{sec:finite-calculus}

For a function on a rational interval $[a,b]$ with a supplied Lipschitz
constant $L$, the left sum on $m$ equal subintervals has error at most
$L(b-a)^2/m$ against any refinement. Use
$m\ge2^{p+4}(1+\ceil{L(b-a)^2})$ and evaluate each summand with
error $2^{-p-4}/(1+b-a)$. The resulting diagonal sequence defines
$\int_a^b f$ as a name.

The finite fundamental theorem follows by telescoping
$F(t_{j+1})-F(t_j)$. If a derivative approximation satisfies
\[
 |F(t+h)-F(t)-h f(t)|\le C h^2,
\]
the sum of its errors is at most $C(b-a)^2/m$. This proves
$\int_a^b f=F(b)-F(a)$ when such a quadratic remainder is supplied.
It applies to our elementary expressions by the finite Taylor series
bounds, and to the gamma expressions by the derivative estimates below.
The product rule with its quadratic remainder proves integration by
parts. Applying this telescoping argument to $F(\phi(t))$, with
explicit bounds for $\phi'$ and its remainder, proves substitution.

Multiple integrals over rectangles are iterated finite sums; changing
their order is a finite rearrangement. We do not need a general
multidimensional change-of-variables theorem: the beta--gamma identity
below is proved by a one-dimensional shift and a finite-product limit.
Singular edges and infinite tails are cut off first using the estimates
in \cref{sec:integrals}. Thus no use is made of Fubini or dominated
convergence for arbitrary functions.

Derivatives of order $r$ mean local polynomial expansions with a
primitive recursive remainder bound. Holomorphic means the corresponding
complex power expansions on every certified disk; entire means that
these bounds are available for every integer radius. Meromorphic means
that on a certified neighborhood a fixed finite power of $(z-z_0)$ times
the function has such an analytic presentation. These are constructive
data, not tests for deciding analyticity of an arbitrary program.

% Mathematical proof material for the main gamma/iPRA article.
% Dependencies: elementary exp and positive-real log with finite error bounds.
% No global complex logarithm, analytic continuation theorem, or choice is used.

\section{The reciprocal function from finite Euler products}
\label{sec:reciprocal-products}

All bounds in this section are uniform in a displayed integer disk radius
$R\geq1$.  Bounds on complex absolute values can be implemented using rational
squared absolute values; an upper bound in the rational norm
$|\Re z|+|\Im z|$ is also sufficient.  Expressions involving nonrational
arguments mean the corresponding fixed primitive recursive names with their
given error certificates.  They do not introduce a quantifier over a universe
of all primitive recursive functions.

\subsection{Euler's constant with a modulus}
Put $H_N=\sum_{k=1}^N1/k$ and $\gamma_N=H_N-\log N$ for $N\geq1$.
The elementary integral or logarithmic-series inequalities give
\[
 0\leq\log(1+1/k)-\frac1{k+1}
 \leq\frac1k-\frac1{k+1}.
\]
Consequently, for $M\geq N$,
\[
 0\leq\gamma_N-\gamma_M\leq\frac1N-\frac1M\leq\frac1N.
 \tag{C1}
\]
Moreover $0\leq\gamma_N\leq1$, by comparing the finite harmonic sum with
the integral of $1/t$ over $[1,N]$.  Thus (C1), together with a finite
approximation of $\log N$, defines a regular name for $\gamma$, and
\[
 0\leq\gamma_N-\gamma\leq N^{-1},\qquad 0\leq\gamma\leq1.
 \tag{C2}
\]
For example $N=2^{p+2}$ and a logarithm evaluation within $2^{-p-2}$
suffice for an approximation with error less than $2^{-p-1}$.
Monotone convergence has not been invoked: (C1) is the Cauchy certificate.

\subsection{A local logarithm that has no branch ambiguity}
For $|w|\leq1/2$ let
\[
 L(w)=\sum_{j=2}^{\infty}\frac{(-1)^{j+1}w^j}{j}.
\]
The series has the explicit estimates
\[
 |L(w)|\leq\frac{|w|^2}{2(1-|w|)}\leq|w|^2,
 \qquad
 \left|L(w)-\sum_{j=2}^{d}\frac{(-1)^{j+1}w^j}{j}\right|
 \leq\frac{|w|^{d+1}}{(d+1)(1-|w|)}.
 \tag{C3}
\]
Here $w+L(w)$ is only the power-series logarithm near $1$.
The identity
\[
 \exp\bigl(w+L(w)\bigr)=1+w
 \tag{C4}
\]
follows directly by finite coefficient algebra.  Write
$F(w)=\exp(w+L(w))=\sum_{n\geq0}f_nw^n$.
The geometric derivative of $w+L(w)$ gives $(1+w)F'=F$ and $f_0=1$.
Comparing finite coefficients yields $f_1=1$ and
$(n+1)f_{n+1}=(1-n)f_n$ for $n\geq1$; hence every $f_n$ with
$n\geq2$ is zero.  Differentiation and multiplication are verified first
for finite polynomials; their geometric and exponential tails bound each
discarded term.  This uses neither a global logarithm on $\mathbb C$
nor an assertion about branches.

\subsection{The entire reciprocal and quantitative Euler convergence}
Define finite entire expressions
\[
 Q_N(z)=z\exp(-z\log N)\prod_{k=1}^N(1+z/k),\qquad
 P_N(z)=z\exp(\gamma z)\prod_{k=1}^N(1+z/k)\exp(-z/k).
 \tag{C5}
\]
Fix $|z|\leq R$ and set $K=2R$.  For $N\geq K$ the tail series
\[
 T_N(z)=\sum_{k=N+1}^{\infty}L(z/k)
\]
has the finite-tail certificate
\[
 |T_N(z)|\leq R^2\sum_{k>N}k^{-2}\leq\frac{R^2}{N}.
 \tag{C6}
\]
The last inequality follows already for every finite sum by
$k^{-2}\leq1/[k(k-1)]$, so that it is a telescoping rational inequality.
Define
\[
 P(z)=z\prod_{k=1}^K(1+z/k)
       \exp\bigl((\gamma-H_K)z+T_K(z)\bigr).
 \tag{C7}
\]
The definitions with different sufficiently large $K$ agree, by (C4) and
finite product cancellation.  In particular
\[
 P(z)=P_N(z)\exp(T_N(z)),\qquad
 Q_N(z)=P_N(z)\exp\bigl(z(\gamma_N-\gamma)\bigr).
 \tag{C8}
\]
The quotients suggested by this notation are never needed at a zero: (C8)
is an equality obtained by multiplication, valid everywhere.

For an entirely rational majorant put
\[
 M_R=R(1+R)^{2R}3^{2R^2+2R},\qquad B_R=R^2+R.
 \tag{C9}
\]
Indeed, the absolute value of the polynomial in (C7) is at most
$R(1+R)^{2R}$, while the exponent has real part at most
\[
 (1+H_{2R})R+R^2/(2R)
 \leq (1+2R)R+R/2\leq2R^2+2R.
\]
Using $e\leq3$ proves $|P(z)|\leq M_R$, and the same estimate bounds
$P_N(z)$ for $N\geq2R$.  By (C2), (C6), (C8), and
$|e^u-e^v|\leq e^{|v|}e^{|u-v|}|u-v|$, or directly by the
exponential power series, one obtains
\[
 |Q_N(z)-P(z)|\leq \frac{3M_RB_R}{N}
 \quad\text{if }N\geq\max(2R,B_R).
 \tag{C10}
\]
One convenient derivation writes both sides using (C7); the difference
of their tail exponents has absolute value at most $B_R/N\leq1$.
Therefore $N\geq3M_RB_R2^p$ gives an error at most $2^{-p}$.
Evaluating the finitely many exponentials and logarithms with a further
error allowance gives a rational algorithm, since all the displayed
lengths and precision budgets are primitive recursive.

Thus (C10) constructs $P$ on every disk.  Equations (C5)--(C8) also prove
the reciprocal Weierstrass product
\[
 \boxed{\displaystyle
 P(z)=z e^{\gamma z}\prod_{k=1}^{\infty}(1+z/k)e^{-z/k}.}
 \tag{C11}
\]
The infinite product means precisely its explicitly controlled finite
products.  It has not been inferred from an abstract product convergence
theorem.

\subsection{An explicit lower bound away from the poles}
Suppose, in addition to $|z|\leq R$, that a rational $0<\delta\leq1$
satisfies
\[
 |z+k|\geq\delta\quad(0\leq k\leq2R).
 \tag{C12}
\]
Set
\[
 \lambda_{R,\delta}=
 \frac{\delta^{2R+1}}{(2R)!\,3^{2R^2+2R}}.
 \tag{C13}
\]
In (C7), the polynomial has absolute value at least
$\delta^{2R+1}/(2R)!$.  The real part of its exponent is at least
$-(2R^2+2R)$, so
\[
 |P(z)|\geq\lambda_{R,\delta}>0.
 \tag{C14}
\]
The identical argument with $\gamma_N$ and the finite tail instead of
$\gamma$ and $T_K$ gives $|Q_N(z)|\geq\lambda_{R,\delta}$ for
$N\geq2R$.  This is a positive lower bound, not merely the negation of an
equality.  It therefore licenses constructive inversion.

On a disk supplied with (C12) define $\Gamma(z)=1/P(z)$.  Its definition
agrees on overlaps.  Its Euler approximants are
\[
 \Gamma_N(z)=\frac{N!N^z}{z(z+1)\cdots(z+N)}=Q_N(z)^{-1}.
\]
For $N\geq\max(2R,B_R)$, (C10)--(C14) yield
\[
 \boxed{\displaystyle
 |\Gamma_N(z)-\Gamma(z)|
 \leq\frac{3M_RB_R}{\lambda_{R,\delta}^{\,2}N}.}
 \tag{C15}
\]
Also $|\Gamma(z)|\leq1/\lambda_{R,\delta}$ and
$|\Gamma(z)|\geq1/M_R$.  In particular gamma has no zeros, with a
quantitative meaning on each supplied disk.  When $\Re z\geq a>0$,
one may take $\delta=\min(a,1)$; hence (C15) identifies this construction
with any independently proved Euler limit on the right half-plane,
including the Euler integral construction.

\subsection{Why ``entire'' is justified constructively}
We give the missing local analytic information rather than applying a
classical theorem about uniformly convergent holomorphic functions.
For $K=2R$, finite rearrangement with the absolute bounds (C3) gives
\[
 T_K(z)=\sum_{m=2}^{\infty}
 \frac{(-1)^{m+1}z^m}{m}
 \left(\zeta(m)-H_K^{(m)}\right),\qquad
 H_K^{(m)}=\sum_{k=1}^K k^{-m}.
 \tag{C16}
\]
Here the positive integer zeta values are defined by their own rational
partial sums, with
\[
 0\leq\zeta(m)-\sum_{k=1}^N k^{-m}
 \leq\frac{1}{(m-1)N^{m-1}}\quad(m\geq2).
 \tag{C17}
\]
The integral comparison in (C17) can equally be replaced by its finite
antiderivative inequality and rational telescoping estimates.
In (C16) the tail after degree $d\geq1$ is at most
\[
 \frac{R^2}{K}\,2^{1-d}.
 \tag{C18}
\]
For example, its $m$th term is bounded by
$(R^2/K)2^{-(m-2)}/[m(m-1)]$, whose geometric tail proves (C18).

Finite polynomial multiplication and the exponential power series now
turn (C7) into a convergent power series on each disk.  This operation has
a direct absolute majorant: if $\sum_j|a_j|t^j\leq A$, then the
coefficient absolute sum of $\exp(\sum_j a_jz^j)$ at $t$ is at most
$e^A$.  This is proved by finite Cauchy multiplication and the ordinary
exponential tail estimate.  There is no appeal to an infinite formal
manipulation without an error bound.

For a useful global consequence, make the construction at radius $2R$
and use $K=4R$.  Its coefficient absolute sum at $2R$ is bounded by
$M_{2R}$.  Consequently, writing $P(z)=\sum_{j\geq1}c_jz^j$,
\[
 \left|P(z)-\sum_{j=1}^d c_jz^j\right|
 \leq M_{2R}2^{-d-1}\quad(|z|\leq R).
 \tag{C19}
\]
The coefficients do not depend on the auxiliary $K$: equality of each
finite-degree coefficient follows by expanding (C4) and cancelling one
factor at a time.  Coefficients are thus computable by finite recursive
convolutions with the constants in (C2) and (C17).

For completeness, a power series about $z_0$, whose coefficient absolute
sum at radius $r$ is at most $A$, satisfies
\[
 |f(z_0+h)-f(z_0)-b_1h|\leq A|h|^2/r^2
 \quad(|h|\leq r),
\]
when $b_1$ is its first coefficient about $z_0$.
The series about any $z_0$ is obtained by the binomial formula, using a
larger disk to bound all rearrangement tails.  The displayed remainder
is the definition and a modulus of its complex derivative.  Repeated
differentiation is handled with an extra disk margin.  This proves
entireness in the concrete sense needed here.  For inversion, suppose
$|P(z_0)|\geq d>0$ and the coefficient absolute sum about $z_0$ is at
most $A$ at radius $r$.  Then $A\geq d>0$.  Choose a positive rational
\[
 \rho\leq\min\left(\frac r2,\frac{dr}{4A}\right).
\]
On $|h|\leq\rho$, the nonconstant coefficient tail gives
$|P(z_0+h)-P(z_0)|\leq A\rho/r\leq d/4$.
The geometric inverse therefore has ratio at most $1/4$ (and, in
particular, at most $1/2$), which provides its convergence and remainder
modulus.  This proves analyticity of $1/P$ on the smaller disk.  Thus
meromorphic language below is supported by local power series and
explicit division bounds.

\section{Functional equation, zeros, poles, and conjugation}

Finite cancellation, without division by $z$, gives
\[
 zQ_N(z+1)=\left(1+\frac{z+1}{N}\right)Q_N(z).
\]
The bounds (C10) give
\[
 P(z)=zP(z+1),\qquad P(1)=1,
 \tag{C20}
\]
the latter also following from $Q_N(1)=(N+1)/N$.  Therefore, wherever
the divisions are licensed,
\[
 \Gamma(z+1)=z\Gamma(z),\qquad\Gamma(1)=1,\qquad
 \Gamma(n+1)=n!.
 \tag{C21}
\]
The last assertion uses ordinary finite induction on $n$; its analytic
content has already been reduced to (C20).

Iteration of (C20) gives the entire identity
\[
 P(z)=\prod_{j=0}^{n}(z+j)\,P(z+n+1).
 \tag{C22}
\]
At $z=-n$ this proves $P(-n)=0$, and differentiating the finite product
proves
\[
 P'(-n)=(-1)^n n!.
 \tag{C23}
\]
So all these zeros are simple.

The converse zero assertion can also be made constructively; it does
not require the invalid general rule that a zero product of arbitrary
constructive reals has a zero factor.  Suppose $P(z)=0$ and $|z|\leq R$.
Compute a rational complex $q$ with $|q-z|\leq1/16$.  Decide the finite
list of rational tests $|q+j|^2<1/16$ for $0\leq j\leq2R$.
If all fail, then $|z+j|\geq3/16$ for all $j$, contradicting (C14).
Hence one test succeeds.  There is at most one such $j$, because two
different integer centers have distance at least $1$.  For this $j$,
$|z+j|<5/16$, while $|z+k|>11/16$ for every other integer $k$.
Factor $(z+j)$ out of (C7).  The remaining factor has a positive lower
bound by exactly the proof of (C14), using $11/16$ for its finitely many
linear factors.  Hence $P(z)=0$ implies $z=-j$.  The bounded search for
$j$ uses a decidable rational predicate and so is primitive recursive.

Equation (C22) gives a local description of gamma at every pole:
\[
 \Gamma(-n+h)=
 \frac{(-1)^n}{n!\,h}\,
 \Gamma(1+h)\prod_{j=1}^n(1-h/j)^{-1}.
 \tag{C24}
\]
All factors apart from $h^{-1}$ have convergent power series for
$|h|<1$.  Their value at zero is $1$, proving a simple pole with residue
\[
 \operatorname{Res}_{z=-n}\Gamma(z)=\frac{(-1)^n}{n!}.
 \tag{C25}
\]
No value is assigned to gamma at its poles.  A totalized implementation
may choose a default value there, but the meromorphic assertions concern
the punctured local formula (C24).

All coefficients and constants in (C5) are real.  Conjugating finite
expressions and using (C10), (C14) proves
\[
 P(\overline z)=\overline{P(z)},\qquad
 \Gamma(\overline z)=\overline{\Gamma(z)}.
 \tag{C26}
\]


\section{Euler's integral and its effective calculus}
\label{sec:integrals}

\subsection{The improper integral is an algorithm}

Fix integers $A,R\ge1$ and a complex input with
\[
 \Rea z\ge1/A,
 \qquad |z|\le R.
\]
For positive $t$ use $t^{z-1}=\exp((z-1)\log t)$. Its modulus is
$t^{\Rea z-1}$, by the finite exponential identities. For
$0<\eta\le1$ and $T\ge1$,
\begin{align}
 \int_0^\eta |t^{z-1}e^{-t}|\dd t
   &\le A\eta^{1/A},\label{eq:gamma-low-tail}\\
 \int_T^\infty |t^{z-1}e^{-t}|\dd t
   &\le 2^{R+1}R!e^{-T/2}.\label{eq:gamma-high-tail}
\end{align}
The first inequality follows by integrating $t^{1/A-1}$. For the
second, $t^{\Rea z-1}\le t^R$ and the single term
$e^{t/2}\ge(t/2)^R/R!$ gives
$t^Re^{-t}\le2^RR!e^{-t/2}$. Initially these arguments apply between
finite positive endpoints; the displayed expressions are the certificates
for removing those endpoints.

Here are actual cutoffs for precision $p$. Put
\[
 D_R=2^{R+1}R!,\quad
 L=p+A+D_R+10,\quad
 \eta=2^{-AL},\quad T=2L.
\]
Then the sum of the two omitted tails is at most
$(A+D_R)2^{-L}<2^{-p-6}$. On $[\eta,T]$ a rational Lipschitz
majorant for the integrand is
\[
 C=(\eta^{-1}+T^R)\bigl((R+1)/\eta+1\bigr).
\]
Indeed its derivative is
$t^{z-1}e^{-t}((z-1)/t-1)$. The finite-calculus mesh of
\cref{sec:finite-calculus}, with $m\ge2^{p+6}(1+CT^2)$,
now specifies every evaluation. This gives a fixed primitive recursive
name $I(z)$ for the Euler integral.

\subsection{Identification with the finite-product gamma function}

For $N\in\N$, integration by parts on $[\eta,1]$ gives
\begin{equation}
 \int_0^1t^{z-1}(1-t)^N\dd t
 =\frac{N!}{z(z+1)\cdots(z+N)}.\label{eq:beta-polynomial}
\end{equation}
In detail, the $N=0$ antiderivative is $t^z/z$ and its lower boundary
has modulus at most $A\eta^{1/A}$. For $N>0$ integration by parts,
with antiderivative $t^z/z$, changes the integral to
$N/z$ times the same integral with $(z,N)$ replaced by $(z+1,N-1)$;
the boundary at $1$ is zero and the lower boundary tends to zero by
the same estimate. Iterating this finite relation proves
\eqref{eq:beta-polynomial}. To read the induction in $\ipra$, carry
the finite mesh and endpoint errors as parameters and decrease each
local error by the product bound \eqref{eq:product-budget}.

The substitution $t=s/N$ gives, for $N\ge1$,
\begin{equation}
 \Gamma_N(z)=\int_0^N s^{z-1}(1-s/N)^N\dd s.
 \label{eq:gamma-euler-integral}
\end{equation}
This formula directly identifies the two constructions with an explicit
rate. If $0\le s\le N/2$, the logarithm series proves
\[
 0\le-N\log(1-s/N)-s\le s^2/N,
\]
and hence
\[
 0\le e^{-s}-(1-s/N)^N\le e^{-s}s^2/N.
\]
On $[N/2,N]$ both kernels are ordered between $0$ and $e^{-s}$;
above $N$ only $e^{-s}$ remains. With
$C_R=1+2^{R+3}(R+2)!$, the elementary monomial bound used for
\eqref{eq:gamma-high-tail} therefore gives
\begin{equation}
 \boxed{\displaystyle
 |\Gamma_N(z)-I(z)|
 \le\frac{C_R}{N}+2^{R+1}R!e^{-N/4}\qquad(N\ge2).}
 \label{eq:integral-product-error}
\end{equation}
For the first term, split its integral at $1$: below $1$,
$s^{\Rea z+1}\le s$, and above $1$ it is at most $s^{R+2}$.
The resulting integral is bounded by $C_R$. The large tail begins at
$N/2$, proving the second term. For an entirely explicit choice set
$N=4(C_R+D_R+p+10)2^{p+4}$. Then each term in
\eqref{eq:integral-product-error} is below $2^{-p-3}$.

Both $I$ and $\Gamma$ are consequently equal to the same effective
Euler limit. We have proved
\begin{theorem}[Euler integral]
For a complex name with a supplied positive lower bound for its real part,
\[
 \boxed{\displaystyle\Gamma(z)=\int_0^\infty t^{z-1}e^{-t}\dd t.}
\]
The integral is absolutely convergent in the finite-tail sense of
\eqref{eq:gamma-low-tail}--\eqref{eq:gamma-high-tail}; it is not an
additional definition requiring an abstract completeness theorem.
\end{theorem}

\subsection{Positivity, modulus, and changes of variable}

For $x>0$, every finite integral summand is nonnegative and the interval
$1\le t\le2$ alone gives
\begin{equation}
 \Gamma(x)\ge1/18>0.\label{eq:gamma-positive}
\end{equation}
Here $t^{x-1}\ge1/2$ and $e^{-t}\ge e^{-2}>1/9$.
On $1/A\le x\le R$ we also have the useful upper bound
\begin{equation}
 \Gamma(x)\le A+2^{R+1}R!.\label{eq:gamma-real-bound}
\end{equation}
The triangle inequality on finite sums, followed by their tails, gives
$|\Gamma(\sigma+iy)|\le\Gamma(\sigma)$ for $\sigma>0$.
There is a stronger vertical-line product formula later in the report.

Integration by parts gives a second proof of
$\Gamma(z+1)=z\Gamma(z)$: the boundary $t^ze^{-t}$ has the same
explicit vanishing estimates. Also $\Gamma(1)=1$ follows from the
antiderivative $-e^{-t}$, so the factorial normalization can be proved
entirely from the integral as well.

Substitution on the truncated interval, followed by the already proved
endpoint bounds, gives
\begin{align}
 \Gamma(z)&=\int_{-\infty}^{\infty}e^{zu-e^u}\dd u,
 \label{eq:exponential-substitution}\\
 \Gamma(z)&=\int_0^1(-\log u)^{z-1}\dd u,
 \label{eq:unit-substitution}\\
 \Gamma(z)&=2c^z\int_0^\infty u^{2z-1}e^{-cu^2}\dd u
 \quad(c>0).\label{eq:quadratic-substitution}
\end{align}
For example \eqref{eq:exponential-substitution} uses
$t=e^u$ and the finite endpoints $u=\log\eta,\log T$;
\eqref{eq:unit-substitution} uses $t=-\log u$ and endpoints
$e^{-T},e^{-\eta}$; \eqref{eq:quadratic-substitution} uses
$t=cu^2$. Nonrational endpoints are evaluated with the local Lipschitz
budget. Thus the substitutions introduce no new convergence principle.
The more general positive-real scaling law is
\begin{equation}
 \int_0^\infty t^{z-1}e^{-ct^r}\dd t
 =\frac{\Gamma(z/r)}{r c^{z/r}}\quad(c,r>0,\ \Rea z>0),
 \label{eq:general-scaling}
\end{equation}
with rational lower and upper bounds for the positive parameters.
It follows from $u=ct^r$, whose derivative is explicit on every
truncated interval. Complex values of $c$ are not included in this
statement; they would require a specified complex power convention and
an additional argument.

\subsection{Differentiation under the integral, with tails}

For every natural number $j$, define the prospective derivative by
\[
 I_j(z)=\int_0^\infty
         t^{z-1}e^{-t}(\log t)^j\dd t.
\]
The additional logarithmic factors have explicit bounds. For
$0<t\le1$,
\[
 |\log t|^j\le j!(2A)^j t^{-1/(2A)},
\]
because $e^{v/(2A)}\ge(v/(2A))^j/j!$ for $v=-\log t$.
For $t\ge1$ use $(\log t)^j\le t^j$. Consequently
\begin{align}
 \int_0^\eta|t^{z-1}e^{-t}(\log t)^j|\dd t
 &\le j!(2A)^{j+1}\eta^{1/(2A)},\label{eq:derivative-low-tail}\\
 \int_T^\infty|t^{z-1}e^{-t}(\log t)^j|\dd t
 &\le2^{R+j+1}(R+j)!e^{-T/2}.\label{eq:derivative-high-tail}
\end{align}
These formulas give primitive recursive cutoff lengths uniformly in $j$.
On the remaining compact interval expand $e^{h\log t}$ to any fixed
degree, using \eqref{eq:exp-tail}. For example the first-order
remainder is at most
\[
 \tfrac12 |h|^2|\log t|^2e^{|h\log t|}.
\]
For $|h|\le1/(4A)$ this is integrable with exactly the same kind of
explicit bounds, using the smaller half-plane margin $1/(2A)$ and
radius $R+1$. Integrating the finite Taylor remainder proves
$I_j'=I_{j+1}$. Thus
\begin{theorem}[All integral derivatives]
For every $j\in\N$ and every certified input in the right half-plane,
\[
 \Gamma^{(j)}(z)
   =\int_0^\infty t^{z-1}e^{-t}(\log t)^j\dd t.
\]
The finite degree and integration errors are controlled by
\eqref{eq:derivative-low-tail}--\eqref{eq:derivative-high-tail} and
the compact Taylor remainder. In particular the statement is not an
unqualified exchange of an infinite limit with differentiation.
\end{theorem}

\subsection{Incomplete gamma as an additional finite-integral consequence}

For $x>0$ and $\Rea z>0$, put
\[
 \gamma_{\ell}(z,x)=\int_0^x t^{z-1}e^{-t}\dd t,
 \qquad \Gamma_u(z,x)=\int_x^\infty t^{z-1}e^{-t}\dd t.
\]
The subscript distinguishes the lower function from Euler's constant.
Additivity of finite sums and integration by parts prove
\begin{align*}
 \gamma_{\ell}(z,x)+\Gamma_u(z,x)&=\Gamma(z),\\
 \gamma_{\ell}(z+1,x)&=z\gamma_{\ell}(z,x)-x^ze^{-x},\\
 \Gamma_u(z+1,x)&=z\Gamma_u(z,x)+x^ze^{-x}.
\end{align*}
For a natural number $n$, induction from $\Gamma_u(1,x)=e^{-x}$ gives
\[
 \boxed{\displaystyle
 \Gamma_u(n+1,x)=n!e^{-x}\sum_{j=0}^n\frac{x^j}{j!}.}
\]
The lower function is $n!$ minus this expression. These formulas extend
to $x=0$ by the explicit small-endpoint estimate, with values $0$ and
$\Gamma(z)$ respectively; no decision whether a general name is zero
is needed in the continuity statement.

% Research contribution: constructive gamma identities.
% This is a section fragment, to be incorporated into the main article.
% All limiting language below abbreviates the displayed finite error estimates.

\section{Finite trigonometry and Euler's reflection formula}

\subsection{The elementary trigonometric foundation}

Here \(\exp\), \(\sin\), and \(\cos\) are the functions given by their
factorial power series, with the explicitly bounded tails described earlier.
Put
\[
 A(t)=\int_0^t\frac{du}{1+u^2},\qquad \pi=4A(1).
\]
The alternating expansion
\[
 A(1)=\sum_{j=0}^{L-1}\frac{(-1)^j}{2j+1}+r_L,
 \qquad |r_L|\le\frac1{2L+1}
\]
is obtained by integrating the finite geometric identity for
\((1+u^2)^{-1}\). Thus \(\pi\) has a primitive recursive rational name.
Finite rational comparisons of the first eight terms give \(3<\pi<4\).
Multiplication of finite exponential series, followed by their factorial
tail bounds, proves the addition laws and
\(\exp(it)=\cos t+i\sin t\).

There is no hidden appeal here to the qualitative theory of trigonometric
functions. For \(0\le t\le1\), differentiate
\[
 e^{-2iA(t)}\frac{1+it}{1-it}.
\]
The two terms cancel exactly. The finite fundamental theorem of calculus
therefore gives
\[
 e^{2iA(t)}=\frac{1+it}{1-it},\qquad
 \sin(2A(t))=\frac{2t}{1+t^2},\qquad
 \cos(2A(t))=\frac{1-t^2}{1+t^2}.                 \tag{T1}
\]
In particular, \(e^{i\pi/2}=i\), \(e^{i\pi}=-1\), and
\(e^{2\pi i}=1\). Since \((2A)'(t)=2/(1+t^2)\ge1\), bounded
bisection constructs its inverse on \([0,\pi/2]\), with an explicit
linear modulus. Formula (T1) proves positivity of sine there. Its second
derivative is \(-\sin\), so the finite second-derivative comparison with
the chord joining the endpoints gives
\[
 \sin v\ge\frac{2v}{\pi}\quad(0\le v\le\pi/2).
                                                        \tag{T2}
\]
For completeness, the chord comparison follows by applying the finite
mean-value inequality to the derivative: a differentiable function with
nonpositive second derivative has decreasing secant slopes. Both this
assertion and (T1) can instead be checked on rational meshes; their
derivative bounds give an explicit mesh error. No unrestricted intermediate
value theorem is used: the only inverse just mentioned has slope at least
one and is obtained by a prescribed bisection program. This program does
not compare exact real values. If the current bracket has width \(w\),
approximate the midpoint value minus the target to error \(w/16\).
When the rational approximation exceeds \(w/8\) in absolute value,
its sign selects a half bracket. Otherwise the slope bound puts the
inverse in the interval of radius \(w/4\) about the midpoint. Intersect
this interval with the old bracket. In every case the width halves, and
every decision is a rational comparison.

\subsection{A finite sine factorization}

Define the entire function
\[
 \operatorname{sinc}(w)=\sum_{j=0}^{\infty}
                \frac{(-1)^jw^{2j}}{(2j+1)!}.
\]
Thus \(\sin w=w\operatorname{sinc}(w)\), including at \(w=0\).
For an odd integer \(m=2J+1\), repeated use of the addition law writes
\(\sin(mu)/(m\sin u)\) as a polynomial of degree \(J\) in
\(\sin^2u\), with constant term one. Its explicitly given distinct roots
are \(\sin^2(k\pi/m)\), \(1\le k\le J\). Consequently the following is
a finite polynomial identity, not an identity obtained by analytic
continuation:
\[
 \operatorname{sinc}(mu)=\operatorname{sinc}(u)
 \prod_{k=1}^{J}
 \left(1-\frac{\sin^2u}{\sin^2(k\pi/m)}\right).           \tag{T3}
\]
To justify the polynomial argument constructively, use the standard
multiple-angle recurrence (equivalently, the recurrence for the
polynomials \(U_j\)). For odd \(m\), \(U_{m-1}(\cos u)\) is a polynomial
in \(1-\sin^2u\). The listed roots are obtained by substituting the
angles \(k\pi/m\), and their distinctness follows from positivity of
cosine on this finite interval. More quantitatively,
\(\cos v\ge1/m\) for
\(0\le v\le(m-1)\pi/(2m)\), by (T2) applied to the complementary
angle. Consecutive sine values therefore differ by at least
\(\pi/m^2\). Repeated monic polynomial division by the known linear
factors proves the factorization; no fundamental theorem of algebra is
needed. Multiplication by \(\operatorname{sinc}(u)\) removes all
division by a possibly zero \(\sin u\).

\subsection{An explicit sine-product error bound}

Let \(R\ge1\) and \(K\ge1\) be integers, let \(|z|\le R\), and put
\[
 m=8K^2R^2+1,\quad J=(m-1)/2,\quad
 S_K(z)=\prod_{k=1}^K\left(1-\frac{z^2}{k^2}\right).
\]
Then
\[
 \left|\operatorname{sinc}(\pi z)-S_K(z)\right|
 \le\frac{2^{74R^2}(R+1)^2}{K}.                         \tag{T4}
\]
Here is a complete estimate. Write
\(u=\pi z/m\), \(v_k=\pi k/m\), and
\(a_k=\sin^2u/\sin^2v_k\). Since
\(\sin v_k\ge2k/m\) and \(|\sin u|\le2|u|\),
\[
 |a_k|\le\frac{16R^2}{k^2},\qquad
 \sum_{k=1}^J|a_k|\le32R^2,
 \qquad \sum_{k=1}^K|z^2/k^2|\le2R^2.                  \tag{T5}
\]
For \(|w|\le1\), the factorial series gives
\(|\operatorname{sinc}(w)-1|\le|w|^2/2\) and
\(|\operatorname{sinc}(w)|\le2\). For \(k\le K\), (T2) also gives
\(\operatorname{sinc}(v_k)>1/2\). Thus
\[
 \left|a_k-\frac{z^2}{k^2}\right|
 \le\frac{96R^2}{m^2}\left(\frac{R^2}{k^2}+1\right).
\]
Indeed, \(|\operatorname{sinc}(u)^2-1|\le3|u|^2/2\), while
\(0<\operatorname{sinc}(v_k)\le1\) implies
\(|1-\operatorname{sinc}(v_k)^2|\le v_k^2\). Dividing their sum by
\(\operatorname{sinc}(v_k)^2>1/4\) bounds the ratio difference by
\(6(|u|^2+v_k^2)\); now use \(\pi^2<16\).
The tail satisfies \(\sum_{k>K}|a_k|\le16R^2/K\). The finite
telescoping identity for a difference of products, together with
\(1+t\le e^t\), therefore yields
\[
 \left|\prod_{k=1}^J(1-a_k)-S_K(z)\right|
 \le e^{34R^2}\left\{\frac{16R^2}{K}
       +\frac{96R^2(2R^2+K)}{m^2}\right\}.              \tag{T6}
\]
By (T3), one must add at most
\(8R^2e^{32R^2}/m^2\) to account for the prefactor
\(\operatorname{sinc}(u)\). Substituting the displayed value of \(m\)
bounds the braces, including this extra term, by \(32R^2/K\).
Finally \(e^{34R^2}\le2^{68R^2}\), which proves the deliberately
generous bound (T4).

This proves, in the precise effective sense needed here,
\[
 \frac{\sin\pi z}{\pi z}
       =\prod_{k=1}^{\infty}\left(1-\frac{z^2}{k^2}\right),             \tag{T7}
\]
where the left side means \(\operatorname{sinc}(\pi z)\) at every
point. For an error \(2^{-p}\), it is enough to take
\(K\ge2^{p+74R^2}(R+1)^2\). Every product and every trigonometric
value occurring in this argument is evaluated to a prescribed rational
precision; its finite error budget is obtained from (T5)--(T6).

\subsection{Reflection and constructive division}

Use the Euler approximants
\[
 \Gamma_N(z)=\frac{N!\,N^z}{z(z+1)\cdots(z+N)}.
\]
At points with the displayed denominator factors apart from zero, exact
finite cancellation gives
\[
 \Gamma_N(z)\Gamma_N(-z)
       =-\frac1{z^2S_N(z)}.                              \tag{R1}
\]
The established convergence of these approximants, the recurrence
\(\Gamma(1-z)=-z\Gamma(-z)\), and (T4) therefore prove
\[
 \boxed{\ \Gamma(z)\Gamma(1-z)=\frac\pi{\sin\pi z}\ }.              \tag{R2}
\]
The domain assertion is constructive: supply a radius \(|z|\le R\)
and a positive rational \(\delta\le1\) separating \(z\) from the
finitely many integers in \([-2R-1,2R+1]\). Integers farther away have
an automatic separation. This is stronger information than the bare
negation ``\(z\) is not an integer,'' and it is exactly what division
requires.

One may extract an explicit lower bound for the sine denominator.
Choose an integer \(L\ge2R\). The first \(L\) factors in (T7) have
product of moduli at least \(\delta^{2L}/(L!)^2\). For the remaining
factors, \(|z^2/k^2|\le1/4\), and the real inequality
\(\log(1-t)\ge-2t\) gives a product lower bound
\(e^{-2R^2/L}\ge2^{-4R^2}\). Hence
\[
 |\sin\pi z|\ge
       \frac{\delta^{2L+1}2^{-4R^2}}{(L!)^2}.             \tag{R3}
\]
This coarse rational bound makes every passage to a reciprocal in (R2)
primitive recursive. Alternatively, the cleared identity
\(\sin(\pi z)\Gamma(z)\Gamma(1-z)=\pi\) can be proved first.

At \(z=1/2\), positivity of the defining gamma integral selects the
positive root and gives
\[
 \boxed{\Gamma(1/2)=\sqrt\pi}.                           \tag{R4}
\]
The square root here is the positive one, computed by bounded rational
bisection. No version of Stirling's formula has been used.

\section{Gauss multiplication without an asymptotic normalization}

\subsection{A second finite sine identity}

For every integer \(m\ge1\),
\[
 \prod_{r=1}^{m-1}\sin\frac{\pi r}{m}
                 =\frac{m}{2^{m-1}}.                    \tag{M1}
\]
For \(m>1\), put \(\omega=e^{2\pi i/m}\). The explicitly known
numbers \(1,\omega,\ldots,\omega^{m-1}\) are pairwise apart: (T2)
gives \(|\omega^r-\omega^s|\ge4/m\) when \(r\ne s\) in this range.
Polynomial division by these known roots proves
\(X^m-1=\prod_{r=0}^{m-1}(X-\omega^r)\). Divide out \(X-1\) and
substitute \(X=1\). Taking squared moduli gives
\[
 m^2=\prod_{r=1}^{m-1}4\sin^2(\pi r/m).
\]
Every sine factor is positive, by (T1), reflection about \(\pi/2\),
and (T2); positive square-root cancellation proves (M1). For \(m=1\)
the empty product is one.

\subsection{Exact grouping of Euler approximants}

Fix \(m\ge1\), set \(a=1/m\), and for \(n\ge1\) set
\(L=m(n+1)-1\). Denominator grouping gives the exact identity
\[
 U_n(z):=\frac{\prod_{r=0}^{m-1}\Gamma_n(z+r/m)}{\Gamma_L(mz)}
   =\frac{(n!)^m m^{m(n+1)}}{L!}\,
       n^{(m-1)/2}\left(\frac nL\right)^{mz}.
\]
All factors independent of \(z\) disappear when this expression is
divided by its value at \(a\):
\[
 \frac{U_n(z)}{U_n(a)}
                 =\left(\frac nL\right)^{m(z-a)}.        \tag{M2}
\]
Only powers of positive real numbers occur, so no branch choice is
implicit. Moreover,
\[
 \left|\log\frac{mn}{L}\right|
   =\log\left(1+\frac{m-1}{mn}\right)\le\frac1n.
\]
For \(|z|\le R\) and \(n\ge m(R+1)\), this proves the explicit
estimate
\[
 \left|\left(\frac nL\right)^{m(z-a)}
           -m^{-m(z-a)}\right|
 \le\frac{2m(R+1)m^{m(R+1)}}n.                           \tag{M3}
\]
Consequently convergence of the gamma approximants proves
\[
 \frac{\prod_{r=0}^{m-1}\Gamma(z+r/m)}{\Gamma(mz)}
   = m^{1-mz} P_m,
 \qquad P_m=\prod_{r=1}^{m-1}\Gamma(r/m).                \tag{M4}
\]
One supplies a positive rational separation of \(mz\) from the
nonpositive integers in the relevant finite disk. This also supplies
separations for all \(z+r/m\): multiplication by \(m\) sends their
possible poles to nonpositive integers. The reciprocal-product bounds
for gamma handle the divisions in (M2)--(M4).

The constant is found without any asymptotic theorem. Reflection and
(M1) give
\[
 P_m^2=\prod_{r=1}^{m-1}
          \Gamma(r/m)\Gamma(1-r/m)
       =\frac{\pi^{m-1}}{\prod_{r=1}^{m-1}\sin(\pi r/m)}
       =\frac{(2\pi)^{m-1}}m.
\]
Since \(P_m>0\), substitution in (M4) proves
\[
 \boxed{\ \prod_{r=0}^{m-1}\Gamma(z+r/m)
       =(2\pi)^{(m-1)/2}m^{1/2-mz}\Gamma(mz)\ }.        \tag{M5}
\]
In particular,
\[
 \Gamma(z)\Gamma(z+1/2)
       =2^{1-2z}\sqrt\pi\,\Gamma(2z),                   \tag{M6}
\]
and
\[
 \Gamma(z)\Gamma(z+1/3)\Gamma(z+2/3)
       =2\pi\,3^{1/2-3z}\Gamma(3z).                    \tag{M7}
\]
Thus duplication is available to determine the constant in Stirling's
formula without circular reasoning.

\section{Beta integrals from one-dimensional finite estimates}

\subsection{Definition, tails, and recurrences}

For certified \(\Re a\ge\alpha>0\), \(\Re b\ge\beta>0\), define
\[
 B(a,b)=\int_0^1t^{a-1}(1-t)^{b-1}\,dt.                 \tag{B1}
\]
All powers use the real logarithm of the positive variable. On
\([\epsilon,1-\epsilon]\), compute a finite rational Riemann sum,
with the explicit derivative bound obtained by differentiating the
integrand. For \(0<\epsilon\le1/2\), the omitted endpoints have
absolute integral at most
\[
 \frac{2\epsilon^\alpha}{\alpha}
       +\frac{2\epsilon^\beta}{\beta}.                  \tag{B2}
\]
This defines a primitive recursive name and a Cauchy modulus; small
rational \(\epsilon=2^{-h}\) is chosen with the displayed bound below
the requested error. Symmetry is the finite substitution \(t\mapsto1-t\).
The identity \(1=t+(1-t)\), and integration of the derivative of
\(t^a(1-t)^b\), give
\[
 B(a,b)=B(a+1,b)+B(a,b+1),\qquad
 aB(a,b+1)=bB(a+1,b).
\]
Indeed, on the truncated interval these are finite integration-by-parts
identities up to the mesh error; the boundary terms have modulus at most
\(\epsilon^\alpha+\epsilon^\beta\). Hence
\[
 B(a+1,b)=\frac a{a+b}B(a,b),\qquad
 B(a,b+1)=\frac b{a+b}B(a,b).                            \tag{B3}
\]
The denominator has real part at least \(\alpha+\beta\), so division
has a supplied margin. Iterating gives
\[
 B(a,b+N)=\frac{(b)_N}{(a+b)_N}B(a,b),                  \tag{B4}
\]
where \((w)_N=\prod_{j=0}^{N-1}(w+j)\). In particular,
\[
 B(a,N+1)=\frac{N!}{a(a+1)\cdots(a+N)}.                \tag{B5}
\]

\subsection{An effective beta-to-gamma limit}

For fixed bounded \(a,b\) as above,
\[
 N^aB(a,b+N)\longrightarrow\Gamma(a).                   \tag{B6}
\]
Here is an estimate sufficient to implement the assertion. Choose
integers \(d\ge\max(1,\Re a)\) and \(C\ge|b-1|\), and write
\(H=1/\alpha+d!\). For \(T\ge1\),
\(N\ge\max(2T,2,T^2+2CT)\), one has
\[
 |N^aB(a,b+N)-\Gamma(a)|
 \le \frac{2H(T^2+2CT)}N
         +2\,4^{d+1}d!\,e^{-T/4}.                     \tag{B7}
\]
To prove it, put \(t=s/N\) in (B1), obtaining an integral on
\([0,N]\). On \([0,T]\), set
\[
 D(s)=(N+b-1)\log(1-s/N)+s.
\]
The elementary inequalities
\(-\log(1-u)-u\le u^2\) and \(|\log(1-u)|\le2u\) for
\(0\le u\le1/2\) yield
\(|D(s)|\le(T^2+2CT)/N\le1\). Therefore
\(|e^{D(s)}-1|\le2|D(s)|\). Since
\(\int_0^\infty s^{\Re a-1}e^{-s}\,ds\le H\), this proves the first
term in (B7). On the remaining interval, \(N+\Re b-1\ge N/2\), so
\[
 |(1-s/N)^{N+b-1}|\le e^{-s/2}.
\]
For \(s\ge T\ge1\), use
\(s^{\Re a-1}\le s^d\) and
\(s^de^{-s/2}\le4^dd!e^{-s/4}\). Integration bounds each of the two
tails by \(4^{d+1}d!e^{-T/4}\), proving (B7).

For example, for requested error \(2^{-p}\), one may choose
\[
 T=4(p+d^2+2d+5),\quad
 N\ge\max\{2T,2,T^2+2CT,
       2^{p+3}(T^2+2CT)\lceil H\rceil\}.
\]
The elementary estimates \(d!\le2^{d^2}\) and \(e\ge2\) verify the
error directly. There is no invocation of dominated convergence.

For \(N\ge2\), the exact Euler-approximant relation is
\[
 N^a\frac{(b)_N}{(a+b)_N}
   =\left(\frac N{N-1}\right)^a
                 \frac{\Gamma_{N-1}(a+b)}{\Gamma_{N-1}(b)}.
\]
Use this in (B4), then apply (B7) and the effective convergence of
\(\Gamma_N\). One obtains first the cleared identity, and then, by the
known nonvanishing margin for gamma,
\[
 \boxed{\ B(a,b)=\frac{\Gamma(a)\Gamma(b)}{\Gamma(a+b)}\ }.          \tag{B8}
\]
This proof uses only one-dimensional substitutions and finite products;
a general change-of-variables theorem for improper double integrals is
unnecessary.

\subsection{Consequences and integral evaluations}

The substitution \(t=u/(1+u)\) in truncated intervals proves
\[
 \int_0^\infty\frac{u^{a-1}}{(1+u)^{a+b}}\,du=B(a,b). \tag{B9}
\]
The omitted tails have bounds \(\epsilon^\alpha/\alpha\) at zero and
\(R^{-\beta}/\beta\) at infinity. The substitution
\(t=\sin^2\theta\), again first on a closed subinterval, gives
\[
 \int_0^{\pi/2}\sin^{2a-1}\theta\cos^{2b-1}\theta\,d\theta
                     =\tfrac12 B(a,b).                \tag{B10}
\]
The beta endpoint estimates (B2) control the removed angular intervals.
For natural numbers \(j,k\),
\[
 B(j+1,k+1)=\frac{j!k!}{(j+k+1)!}.                      \tag{B11}
\]
Combining (B8)--(B9) with reflection gives, for certified
\(0<\Re a<1\),
\[
 \int_0^\infty\frac{t^{a-1}}{1+t}\,dt
       =\frac\pi{\sin\pi a}.                          \tag{B12}
\]
For positive rational \(q\) and \(0<\Re a<q\), replacing \(t\) by
\(u^{1/q}\) gives
\[
 \int_0^\infty\frac{t^{a-1}}{1+t^q}\,dt
       =\frac\pi{q\sin(\pi a/q)}.                     \tag{B13}
\]
Similarly, for positive rationals \(c,q\),
\[
 \int_0^\infty t^{a-1}e^{-ct^q}\,dt
       =\frac{\Gamma(a/q)}{q\,c^{a/q}},\qquad \Re a>0. \tag{B14}
\]
All substitutions are between explicitly truncated intervals; their
endpoint errors follow from (B2), (B7), or the gamma tail estimates.

\section{Special values and exact rational-argument relations}

Recurrence and (R4) give, for \(n\ge0\),
\[
 \Gamma(n+1/2)=\frac{(2n)!}{4^n n!}\sqrt\pi,\qquad
 \Gamma(1/2-n)=\frac{(-4)^n n!}{(2n)!}\sqrt\pi.          \tag{V1}
\]
The finite verification is induction: each additional factor is
\(n+1/2\) in the first formula and \(1/2-n\) in the second.
Thus, for example, \(\Gamma(-1/2)=-2\sqrt\pi\) and
\(\Gamma(-3/2)=4\sqrt\pi/3\).

Reflection supplies
\[
 \Gamma(1/3)\Gamma(2/3)=\frac{2\pi}{\sqrt3},\quad
 \Gamma(1/4)\Gamma(3/4)=\pi\sqrt2,\quad
 \Gamma(1/6)\Gamma(5/6)=2\pi.                           \tag{V2}
\]
Let \(A=\Gamma(1/3)\), \(D=\Gamma(1/4)\), both positive. Duplication
at \(1/6\), followed by (V2), gives
\[
 \Gamma(1/6)=\frac{2^{-1/3}\sqrt3}{\sqrt\pi}A^2,
 \qquad
 \Gamma(5/6)=\frac{2^{4/3}\pi^{3/2}}{\sqrt3\,A^2}.     \tag{V3}
\]
Duplication and triplication also give the useful twelfth-argument
relations
\[
 \Gamma(1/12)\Gamma(7/12)=\sqrt6\,A^2,
 \qquad
 \Gamma(1/12)\Gamma(5/12)=\sqrt2\,3^{1/4}D^2.           \tag{V4}
\]
Combining these with reflection proves the explicit positive-square
identities
\[
 \Gamma(1/12)^2
   =\frac{3^{3/4}(\sqrt6+\sqrt2)}{2\pi}A^2D^2,
 \qquad
 \Gamma(5/12)^2
   =\frac{\pi(\sqrt6-\sqrt2)}{3^{1/4}}\frac{D^2}{A^2}. \tag{V5}
\]
All roots in these expressions are positive roots. They are computed by
bounded bisection, and their algebraic identities are verified by
polynomial multiplication together with positive lower bounds.

For fifths, reflection yields
\[
 \Gamma(1/5)\Gamma(4/5)
       =\frac{4\pi}{\sqrt{10-2\sqrt5}},\qquad
 \Gamma(2/5)\Gamma(3/5)
       =\frac{4\pi}{\sqrt{10+2\sqrt5}}.                 \tag{V6}
\]
More generally, the finite normalization computation already proved
\[
 \prod_{r=1}^{m-1}\Gamma(r/m)
          =\frac{(2\pi)^{(m-1)/2}}{\sqrt m}.             \tag{V7}
\]
These are relations among values, not claims that all rational-argument
gamma values reduce to elementary constants. A finite certificate for a
proposed relation can consist of applications of recurrence, reflection,
and multiplication, followed by exact rational and algebraic arithmetic.

For the rising factorial, the recurrence also proves
\[
 \Gamma(z+n)=(z)_n\Gamma(z),                            \tag{V8}
\]
at arguments with the requisite pole margins. For integer \(n\ge0\),
the generalized binomial polynomial satisfies
\[
 \binom zn=\frac{\Gamma(z+1)}{n!\,\Gamma(z-n+1)}        \tag{V9}
\]
whenever the displayed gamma values are defined. The left side is the
finite polynomial \(z(z-1)\cdots(z-n+1)/n!\). At an apparent cancelled
pole, use this polynomial or a reciprocal-gamma version, rather than
performing undefined division.

\section{Modulus formulas on vertical lines}

Complex conjugation commutes with every Euler approximant and hence
with gamma. Apply reflection to \(z=iy\), with real \(y\) apart from
zero, and use \(\sin(i\pi y)=i\sinh(\pi y)\). Recurrence gives
\[
 |\Gamma(iy)|^2=\frac\pi{y\sinh(\pi y)},\qquad
 |\Gamma(1+iy)|^2=\frac{\pi y}{\sinh(\pi y)}.           \tag{VM1}
\]
For all real \(y\), including zero, a uniformly constructive version is
\[
 |\Gamma(1+iy)|^2
  =\left(\sum_{j=0}^{\infty}
                \frac{(\pi y)^{2j}}{(2j+1)!}\right)^{-1}.             \tag{VM2}
\]
The denominator is at least one. To verify (VM2) without cancelling
\(y\), use the exact finite identity
\[
 \Gamma_N(1+iy)\Gamma_N(1-iy)
   =\frac{N^2}{(N+1)^2}
     \left(\prod_{k=1}^{N+1}(1+y^2/k^2)\right)^{-1}.
\]
The entire sine product at \(z=iy\) identifies the limit of this
positive product with the denominator in (VM2). Every reciprocal has
lower bound one before inversion, so no decision whether \(y=0\) is
required.
At \(z=1/2+iy\), reflection instead gives
\[
 |\Gamma(1/2+iy)|^2=\frac\pi{\cosh(\pi y)}              \tag{VM3}
\]
for every real \(y\); this denominator too is bounded below by one.
Repeated recurrence yields
\[
 |\Gamma(n+iy)|^2
  =\prod_{k=1}^{n-1}(k^2+y^2)\,|\Gamma(1+iy)|^2
       \quad(n\ge1),                                  \tag{VM4}
\]
and
\[
 |\Gamma(n+1/2+iy)|^2
  =\prod_{k=0}^{n-1}((k+1/2)^2+y^2)
                   \frac\pi{\cosh(\pi y)}\quad(n\ge0). \tag{VM5}
\]

\section{Differentiated functional identities}

The derivative estimates for the reciprocal product justify
differentiating the finite identities above with their controlled
remainders. With \(\psi=\Gamma'/\Gamma\), one obtains
\[
 \psi(1-z)-\psi(z)=\pi\cot(\pi z),\qquad
 \psi'(1-z)+\psi'(z)=\pi^2\csc^2(\pi z),                \tag{D1}
\]
and
\[
 \sum_{r=0}^{m-1}\psi(z+r/m)=m\psi(mz)-m\log m.         \tag{D2}
\]
For \(j\ge1\), further differentiation gives
\[
 \sum_{r=0}^{m-1}\psi^{(j)}(z+r/m)
                         =m^{j+1}\psi^{(j)}(mz).        \tag{D3}
\]
Combining these with \(\psi(1)=-\gamma\) yields
\[
 \begin{aligned}
 \psi(1/2)&=-\gamma-2\log2,\\
 \psi(1/3)&=-\gamma-\frac{\pi}{2\sqrt3}-\frac32\log3,\quad
 \psi(2/3)=-\gamma+\frac{\pi}{2\sqrt3}-\frac32\log3,\\
 \psi(1/4)&=-\gamma-\frac\pi2-3\log2,\quad
 \psi(3/4)=-\gamma+\frac\pi2-3\log2,\\
 \psi(1/6)&=-\gamma-\frac{\pi\sqrt3}{2}-\frac32\log3-2\log2,\\
 \psi(5/6)&=-\gamma+\frac{\pi\sqrt3}{2}-\frac32\log3-2\log2.
 \end{aligned}                                        \tag{D4}
\]
For example, reflection at \(1/3\) determines the difference of the
two third-argument values, while (D2) with \(m=3,z=1/3\) determines
their sum. The quarter- and sixth-argument formulas follow in exactly
the same finite way from duplication and reflection. Thus (D4) involves
only solving two linear equations with nonzero integer coefficients.
At \(z=1/2\), (D1) gives \(\psi'(1/2)=\pi^2/2\).

\subsection{How the identities are read in iPRA}

For each fixed collection of primitive recursive input names and
certificates, the identities in this section mean equality of their
rational approximants at every prescribed accuracy. Bounds (T4), (R3),
(M3), (B2), and (B7), together with the gamma-product bounds, determine
all truncation lengths and arithmetic precisions by primitive recursive
formulas. Product regrouping, polynomial division by the explicitly
given roots, integration by parts on a finite mesh, and all endpoint
comparisons have finite coded invariants. Quantifier-free induction
proves those invariants. The language does not quantify over an
uncertified universal evaluator for arbitrary primitive recursive
programs. The claims are uniform finite arithmetic lemmas and their
instances for the stated names, rather than an implicit appeal to a
higher-order theory of all continuous functions.

% Primary formula references, checked 2026-11-15:
% NIST DLMF 5.5, https://dlmf.nist.gov/5.5 : recurrence, reflection,
% duplication, multiplication. Retrieval reference turn3view0.
% NIST DLMF 5.12, https://dlmf.nist.gov/5.12 : beta integrals.
% Retrieval reference turn3view1.
% The finite proofs and intentionally coarse moduli above are supplied
% here independently; citations identify the standard identities only.

\section{Logarithmic derivatives and explicit series}

For an input equipped with pole separation define
$\psi(z)=\Gamma'(z)/\Gamma(z)$; this division is licensed by
$|\Gamma(z)|\geq1/M_R$.  Differentiating the finite factorization (C7)
and its explicitly convergent power series gives
\[
 \boxed{\displaystyle
 \psi(z)=-\gamma-\frac1z+
       \sum_{k=1}^{\infty}\frac{z}{k(k+z)}.}
 \tag{C27}
\]
For $N\geq2R$ the tail after $N$ is at most $2R/N$.
Indeed $|k+z|\geq k/2$, and therefore the $k$th term is at most
$2R/k^2$.  This also explains why the separated form is useful: each
finite head term has a certified denominator, and the tail has a
uniform bound independent of the fine pole separation.

An equivalent partial fraction series is
\[
 \psi(z)=-\gamma+
 \sum_{k=0}^{\infty}\left(\frac1{k+1}-\frac1{k+z}\right).
 \tag{C28}
\]
The difference of corresponding finite sums is an explicitly vanishing
endpoint term.  For every integer $m\geq1$,
\[
 \boxed{\displaystyle
 \psi^{(m)}(z)=(-1)^{m+1}m!
               \sum_{k=0}^{\infty}(k+z)^{-m-1}.}
 \tag{C29}
\]
If $N\geq2R$, the tail over $k>N$ has absolute value at most
\[
 \frac{(m-1)!2^{m+1}}{N^m}.
 \tag{C30}
\]
These differentiations are not justified merely by convergence of
(C27).  On a slightly larger pole-separated disk, differentiating
$L(z/k)$ $r$ times gives, for $r\geq2$,
\[
 \frac{d^r}{dz^r}L(z/k)
   =\frac{(-1)^{r-1}(r-1)!}{(k+z)^r},
\]
with tail bounded by
$(r-1)!2^r/[(r-1)N^{r-1}]$.
For $r=1$ its tail is at most $2R/N$.
Finite Taylor remainders, bounded by one further derivative on that
larger disk, prove the derivative assertions with explicit moduli.
All lengths, factorials, and tail budgets remain primitive recursive
uniformly in the requested derivative order.

Differentiating the functional equation or reindexing finite sums gives
\[
 \psi(z+1)=\psi(z)+z^{-1},\qquad
 \psi^{(m)}(z+1)=\psi^{(m)}(z)+(-1)^m m!z^{-m-1}.
 \tag{C31}
\]
In particular, for $n\geq0$ and $m\geq1$,
\[
 \begin{aligned}
 \psi(n+1)&=H_n-\gamma,\\
 \psi^{(m)}(1)&=(-1)^{m+1}m!\zeta(m+1),\\
 \psi^{(m)}(n+1)&=(-1)^{m+1}m!
       \bigl(\zeta(m+1)-H_n^{(m+1)}\bigr).
 \end{aligned}
 \tag{C32}
\]
Separation of even and odd terms in (C29), a finite reindexing followed
by (C30), yields
\[
 \psi^{(m)}(1/2)=(-1)^{m+1}m!(2^{m+1}-1)\zeta(m+1).
 \tag{C33}
\]
Similarly the $N$th partial sum in (C28) at $z=1/2$ is
$-\gamma+2H_N-2H_{2N}$.  By (C2),
$H_{2N}-H_N-\log2=\gamma_{2N}-\gamma_N$ has absolute value at most
$1/N$.  Hence
\[
 \psi(1/2)=-\gamma-2\log2.
 \tag{C34}
\]
For positive real $x$ the series in (C29) has positive terms, giving
the quantitative sign statement
\[
 (-1)^{m+1}\psi^{(m)}(x)\geq m!x^{-m-1}>0.
 \tag{C35}
\]
In particular $\psi$ is strictly increasing and $\log\Gamma$ is
strictly convex on the positive real axis.  On a rational interval
$[a,b]$ the lower bound $1/b^2$ for $\psi'$ supplies a modulus for
strict convexity; it is stronger than a purely qualitative assertion.

Higher gamma derivatives can be computed without new analytic arguments:
\[
 \Gamma^{(n)}(z)=\Gamma(z)Y_n(z),\quad Y_0=1,\qquad
 Y_{n+1}=\sum_{j=0}^n\binom nj Y_{n-j}\psi^{(j)}.
 \tag{C36}
\]
This is the ordinary finite Leibniz rule applied to
$\Gamma'=\Gamma\psi$, and is a primitive recursive polynomial
recursion for every derivative order.

\section{Taylor and Laurent expansions with remainders}

For $|w|\leq r<1$, (C11), (C20), and finite geometric bounds prove
\[
 \Gamma(1+w)=\exp\!\left(
 -\gamma w+\sum_{m=2}^{\infty}
                    \frac{(-1)^m\zeta(m)}m w^m\right).
 \tag{C37}
\]
The cancellation of the factor $w$ is cancellation of power-series
coefficients (or of the corresponding finite polynomial factors),
so this argument also applies at $w=0$ without deciding whether an
arbitrary named complex number is zero.
The exponent is the local branch of $\log\Gamma(1+w)$ whose value at
$w=0$ is zero.  The wording is essential: no global principal logarithm
identity is asserted.  Since $1\leq\zeta(m)\leq2$ for $m\geq2$,
the tail after degree $d\geq1$ of this exponent is at most
\[
 \frac{2r^{d+1}}{(d+1)(1-r)}.
 \tag{C38}
\]
The same argument, or (C27), gives
\[
 \psi(1+w)=-\gamma+
       \sum_{m=1}^{\infty}(-1)^{m+1}\zeta(m+1)w^m,
 \tag{C39}
\]
whose tail after degree $d$ is at most $2r^{d+1}/(1-r)$.
To turn these into explicit primitive recursive degree bounds, choose
an integer $q\geq2$ with $r\leq1-1/q$.  The finite binomial inequality
gives $r^q\leq1/2$, so a degree of the form $q(p+c)$ reduces the
geometric factor below $2^{-p-c}$.  The remaining prefactors are
absorbed by an explicitly chosen integer $c$.

Writing $\Gamma(1+w)=\sum_{n\geq0}g_nw^n$, finite coefficient
multiplication gives
\[
 g_0=1,\qquad
 ng_n=-\gamma g_{n-1}+
         \sum_{m=2}^n(-1)^m\zeta(m)g_{n-m}.
 \tag{C40}
\]
For example
\[
 g_1=-\gamma,\quad
 g_2=\frac{\gamma^2+\zeta(2)}2,\quad
 g_3=-\frac{\gamma^3+3\gamma\zeta(2)+2\zeta(3)}6.
\]
Thus $\Gamma'(1)=-\gamma$ and
$\Gamma''(1)=\gamma^2+\zeta(2)$.
If a direct Taylor-polynomial error is desired on $|w|\leq r<1$,
choose a rational $r<s<1$.  The coefficient absolute sum at $s$ is at
most
\[
 A_s=\exp\left(s+\frac{s^2}{1-s}\right),
\]
because $2/m\leq1$ for $m\geq2$.  Hence
\[
 \left|\Gamma(1+w)-\sum_{n=0}^d g_nw^n\right|
 \leq A_s(r/s)^{d+1}.
 \tag{C41}
\]
An integer power of $3$ can replace $A_s$ to make the majorant rational.

Writing $P(z)=\sum_{n\geq1}c_nz^n$, one obtains
\[
 c_1=1,\quad c_2=\gamma,\qquad
 (n-1)c_n=\gamma c_{n-1}+
    \sum_{m=2}^{n-1}(-1)^{m+1}\zeta(m)c_{n-m}\quad(n\geq3).
 \tag{C42}
\]
For example
\[
 c_3=\frac{\gamma^2-\zeta(2)}2,\qquad
 c_4=\frac{\gamma^3-3\gamma\zeta(2)+2\zeta(3)}6.
\]
Although its derivation starts in $|z|<1$, this reciprocal Maclaurin
series converges on every disk by (C19), with the same coefficients.

There is also a concrete Laurent remainder at every pole.  For
$0<|h|\leq1/4$, combine (C24) and (C37) to write
\[
 \Gamma(-n+h)=\frac{(-1)^n}{n!h}\exp\bigl(Ah+V(h)\bigr),
 \quad A=H_n-\gamma,
\]
where
\[
 V(h)=\sum_{m=2}^{\infty}
 \frac{(-1)^m\zeta(m)+H_n^{(m)}}m h^m,
 \qquad |V(h)|\leq3|h|^2.
\]
Since $|A|\leq n+1$, the exponential remainder gives
\[
 \left|\Gamma(-n+h)-\frac{(-1)^n}{n!}
       \left(\frac1h+H_n-\gamma\right)\right|
 \leq\frac{C_n}{n!}|h|,
 \quad C_n=3+3^{n+2}(n+2)^2.
 \tag{C43}
\]
The generous integer constant avoids any issue of computing an
exponential bound.  At $n=0$ this includes
$\Gamma(h)=h^{-1}-\gamma+O(h)$ with the displayed effective meaning.

\subsection{Taylor series at every positive integer}
The same cancellation gives a wider family of branch-free exponential
formulas.  For $n\geq1$ and $|w|<n$,
\[
 \Gamma(n+w)=(n-1)!\exp\!\left((H_{n-1}-\gamma)w+
   \sum_{m=2}^{\infty}\frac{(-1)^m}{m}
          \bigl(\zeta(m)-H_{n-1}^{(m)}\bigr)w^m\right).
 \tag{C44}
\]
This follows either by cancelling the first $n-1$ factors of (C11),
or by multiplying (C37) by $\prod_{j=1}^{n-1}(j+w)$ and comparing
finite coefficients.  The former proof immediately supplies the full
disk $|w|<n$, without an analytic continuation theorem.  Since
\[
 0<\zeta(m)-H_{n-1}^{(m)}
       \leq n^{-m}+\frac{n^{1-m}}{m-1}\leq(n+1)n^{-m},
\]
the exponent tail after degree $d\geq1$, on $|w|\leq r<n$, is at most
\[
 \frac{(n+1)(r/n)^{d+1}}{(d+1)(1-r/n)}.
 \tag{C45}
\]
For $n=2$ this is the branch-free version of the commonly quoted
log-gamma expansion involving $\zeta(m)-1$.  If it is rewritten as a
formula for $\log\Gamma(1+w)$, the additional logarithm of $1+w$
requires a compatible local branch; (C44) avoids that issue entirely.

\section{Two further finite-product consequences}

Let $x,y$ be real and supply pole separation for $x$.  Taking absolute
values in the finite Euler quotient gives the exact finite identity
\[
 \left|\frac{\Gamma_N(x)}{\Gamma_N(x+iy)}\right|^2
       =\prod_{k=0}^N\left(1+\frac{y^2}{(x+k)^2}\right).
\]
Here $|N^{iy}|=1$, an identity of the elementary exponential.
The left side has a specified limit by (C15).  Alternatively, the
logarithm of the right-hand tail is bounded by $4y^2/N$ as soon as
$N\geq2|x|$, since $\log(1+t)\leq t$ for $t\geq0$.
Thus
\[
 \boxed{\displaystyle
 \left|\frac{\Gamma(x)}{\Gamma(x+iy)}\right|^2
       =\prod_{k=0}^{\infty}\left(1+\frac{y^2}{(x+k)^2}\right).}
 \tag{C46}
\]
In particular $|\Gamma(x+iy)|\leq|\Gamma(x)|$, and if a positive
lower bound for $|y|$ is supplied, the inequality is strict with a
quantitative margin already furnished by the $k=0$ factor.  No
positivity assumption on $x$ is needed, only its separation from the
nonpositive integers.

More generally, let finite lists $a_1,\ldots,a_m$ and
$b_1,\ldots,b_m$ satisfy $\sum_i a_i=\sum_i b_i$.  Supply pole
separation for the $b_i$; the $a_i$ need no such condition.  Finite
cancellation proves
\[
 \prod_{k=0}^{N}\prod_{i=1}^m\frac{a_i+k}{b_i+k}
       =\prod_{i=1}^m\frac{Q_N(a_i)}{Q_N(b_i)}.
\]
Hence
\[
 \boxed{\displaystyle
 \prod_{k=0}^{\infty}\prod_{i=1}^m\frac{a_i+k}{b_i+k}
       =\prod_{i=1}^m P(a_i)\Gamma(b_i).}
 \tag{C47}
\]
If the $a_i$ are also separated from the poles, the right side is
$\prod_i\Gamma(b_i)/\Gamma(a_i)$.  Allowing $P(a_i)$ instead makes
the zero cases meaningful as well.  Suppose all the parameters have
absolute value at most $R$ and all $b_i$ have separation $\delta$.
The proof of (C14) also bounds $|Q_N|$ by $M_R$.  The difference of
one quotient in the last finite display and its limit is therefore
at most $6M_R^2B_R/(\lambda_{R,\delta}^2N)$.
The finite telescoping product identity yields the explicit total error
\[
 \frac{6mM_R^2B_R}{\lambda_{R,\delta}^2N}
       \left(\frac{M_R}{\lambda_{R,\delta}}\right)^{m-1}
 \quad(m\geq1).
 \tag{C48}
\]
The empty-list case is the exact equality $1=1$.

\paragraph{Reference alignment and formalization boundary.}
The identities in this development agree with NIST DLMF
\S5.2 (definitions and poles), \S5.7 (series), \S5.8 (products), and
\S5.15 (polygamma functions).  The proof above supplies its own finite
estimates.  Every infinite statement is read as its family of rational
error inequalities.  A Lean theorem stated directly over Mathlib's
$\mathbb C$ is a useful presentation theorem but is not, by itself,
an iPRA proof.  The arithmetic layer must verify the finite products,
tail bounds, coefficient recursions, and precision transformers, with
quantifier-free induction on their displayed finite recursions.

% LaTeX fragment.  Uses standard theorem/lemma/proof environments.
% Prepared independently, with the finite telescoper in the repository's
% Stirling draft checked and extended from integers to positive real names.
\section{Real convexity, effective uniqueness, and gamma inequalities}
\label{sec:real-convexity}

Throughout this section an argument $x>0$ is supplied with rational bounds
$0<a\le x\le b$.  Every real-valued expression denotes its fixed primitive
recursive rational approximation procedure.  An assertion with arbitrary
such procedures is a theorem schema, not quantification over an internally
available universe of real functions.  Integral notation below abbreviates
finite rational quadrature with the indicated error estimates.

\subsection{Digamma estimates and strong log-convexity}

Put $g(x)=\log\Gamma(x)$ and $\psi(x)=g'(x)$.  Differentiating the finite
Euler approximants gives
\[
 g_N'(x)=\log N-\sum_{k=0}^{N}\frac1{x+k},\qquad
 g_N''(x)=\sum_{k=0}^{N}\frac1{(x+k)^2}.
\]
The product construction supplies their uniform convergence on $[a,b]$.
For the second derivatives a sufficient tail estimate is
\[
 0\le\sum_{k=N+1}^{M}\frac1{(x+k)^2}
 \le\int_N^M\frac{dt}{(x+t)^2}\le\frac1{x+N}.
\]
Consequently
\begin{equation}\label{eq:psi-prime-positive}
 \psi'(x)=\sum_{k=0}^{\infty}\frac1{(x+k)^2},\qquad
 \frac1{b^2}\le\psi'(x)\le\frac1{a^2}+\frac1a
 \quad(a\le x\le b).
\end{equation}
To justify differentiation without a general interchange theorem, apply
the finite fundamental theorem to $g_N'$ on a rational segment, replace its
second derivative by a partial sum, and use the displayed uniform tail.
The resulting finite increment identities have an error tending to zero
at the displayed primitive recursive rate.  The same argument for $g_N$
identifies its first derivative.

For completeness, comparison of $t\mapsto(x+t)^{-1}$ with its upper and
lower rectangle sums gives
\[
 \log x-\frac1x\le\psi(x)\le\log x.
\]
The upper bound can be sharpened.  Convexity of this reciprocal function
gives, on each unit interval, that its integral is at most its trapezoidal
sum.  Summing and substituting in the formula for $g_N'$ yields
\begin{equation}\label{eq:digamma-elementary-bounds}
 \log x-\frac1x\le\psi(x)\le\log x-\frac1{2x}.
\end{equation}
All discarded terminal terms are bounded by rational multiples of $1/N$;
for example $0\le\log(1+x/N)\le b/N$.

Let $a\le x,y\le b$ and $0\le s\le1$.  The function
$g(t)-t^2/(2b^2)$ has nonnegative second derivative.  Integrating its
nondecreasing first derivative on the two sides of a division point proves
the secant inequality, giving
\begin{equation}\label{eq:strong-log-convexity}
 g((1-s)x+sy)
 \le(1-s)g(x)+sg(y)-\frac{s(1-s)(x-y)^2}{2b^2}.
\end{equation}
One can derive the same conclusion directly by adding finite interval
increments; no mean-value point is chosen.  Exponentiating proves
log-convexity of gamma.  The explicit last term proves strictness whenever
$s$, $1-s$, and $|x-y|$ have positive rational lower bounds.

In particular gamma is strictly convex as well: differentiating once more
gives $\Gamma''=\Gamma(\psi^2+\psi')>0$.  This conclusion has an explicit
positive gap on each interval $[a,b]$, using a positive lower bound for
$\Gamma$ there and \eqref{eq:psi-prime-positive}.

\subsection{A Bohr--Mollerup squeeze which needs no floor of a real}

Let $F$ be a fixed positive primitive recursive presentation satisfying
$F(1)=1$, $F(x+1)=xF(x)$, and log-convexity, with local positive bounds so
that $f=\log F$ is defined.  Choose an integer $B\ge\max(1,b)$ and an
integer $n\ge2$.  Convexity, first using the points $n-1,n,n+x$, and then
the points $n,n+x,n+B$, gives
\[
 f(n)+x\log(n-1)
 \le f(n+x)
 \le f(n)+\frac{x}{B}\sum_{j=0}^{B-1}\log(n+j)
 \le f(n)+x\log(n+B-1).
\]
Here $f(n)=\log((n-1)!)$, and the recurrence gives
$f(n+x)=f(x)+\sum_{k=0}^{n-1}\log(x+k)$.  Therefore $f(x)$ belongs to
the explicit interval
\begin{align*}
 L_n(x)&=\log((n-1)!)+x\log(n-1)
              -\sum_{k=0}^{n-1}\log(x+k),\\
 U_n(x)&=\log((n-1)!)+x\log(n+B-1)
              -\sum_{k=0}^{n-1}\log(x+k).
\end{align*}
Its width satisfies
\begin{equation}\label{eq:bm-width-general}
 0\le U_n(x)-L_n(x)
 =x\log\left(1+\frac{B}{n-1}\right)
 \le\frac{B^2}{n-1}.
\end{equation}
The already constructed gamma presentation satisfies the same hypotheses,
so both $f(x)$ and $g(x)$ lie in this interval.  Choosing
$n=1+B^2 2^{p+2}$ proves $|f(x)-g(x)|\le2^{-p-2}$, apart from elementary
function approximation errors which are assigned the remaining budget.
Thus $F(x)=\Gamma(x)$ as equality of names.  Exponentiation is legitimate
with the explicit local bound supplied either in the presentation of $F$
or by the displayed squeeze at $n=2$.

This proof uses only two convexity instances and $n+B$ recurrence instances
at each requested precision.  It does not find the integer part of an
arbitrary real, nor use a completeness principle for arbitrary functions.
For rational arguments, log-convexity may equivalently be supplied by the
integer-power inequality
$F((u x+v y)/(u+v))^{u+v}\le F(x)^uF(y)^v$.

\subsection{Wendel, Gautschi, and effective gamma-ratio asymptotics}

For $x>0$ and $0\le s\le1$, log-convexity between $x$ and $x+1$ gives
$\Gamma(x+s)\le x^s\Gamma(x)$.  Applying the same inequality to $x+s$
with increment $1-s$ gives
$x\Gamma(x)\le(x+s)^{1-s}\Gamma(x+s)$.  Thus
\begin{equation}\label{eq:wendel-constructive}
 \left(\frac{x}{x+s}\right)^{1-s}
 \le\frac{\Gamma(x+s)}{x^s\Gamma(x)}\le1.
\end{equation}
For $0<s<1$ both inequalities are strict, with gaps obtainable from
\eqref{eq:strong-log-convexity}.  This also gives Gautschi's inequalities
\begin{equation}\label{eq:gautschi-constructive}
 x^{1-s}<\frac{\Gamma(x+1)}{\Gamma(x+s)}<(x+1)^{1-s}
 \qquad(0<s<1).
\end{equation}
Indeed the first follows from the preceding upper bound for
$\Gamma(x+s)$; the preceding second application gives the stronger upper
bound $(x+s)^{1-s}<(x+1)^{1-s}$.

Writing the middle expression of \eqref{eq:wendel-constructive} as $Q$,
\[
 -\frac{s(1-s)}x\le\log Q\le0,
 \qquad 0\le1-Q\le\frac{s(1-s)}x\le\frac1{4x}.
\]
This is an explicit uniform gamma-ratio asymptotic, not just a limit
assertion.  More generally, let $|\alpha|,|\beta|\le A$ with $A\ge1$ and
$x\ge2A+1$.  On the segment joining $x+\alpha$ to $x+\beta$,
\eqref{eq:digamma-elementary-bounds} implies
\[
 |\psi(x+t)-\log x|
 \le\frac{2(A+1)}x\quad(|t|\le A).
\]
Integration over this finite segment gives
\begin{equation}\label{eq:general-gamma-ratio}
 \left|\log\frac{\Gamma(x+\alpha)}{\Gamma(x+\beta)x^{\alpha-\beta}}\right|
 \le\frac{4A(A+1)}x.
\end{equation}
If $C=4A(A+1)$ and $x\ge\max(2A+1,C)$, the relative error is at most
$2C/x$, since $|e^u-1|\le2|u|$ for $|u|\le1$.

\section{Critical points, including every negative interval}
\label{sec:gamma-critical-points}

The standard harmonic approximation to Euler's constant satisfies
\[
 H_N-\log(N+1)\le\gamma\le H_N-\log N,
\]
with bracket width at most $1/N$.  The elementary series
$\log2=2\sum_{j\ge0}3^{-2j-1}/(2j+1)$ gives
$2/3<\log2<3/4$.  Taking $N=1$ in the lower bracket and $N=4$ in the
upper bracket proves
\[
 \frac14<\gamma<\frac34.
\]
The digamma series therefore gives
$\psi(1)=-\gamma<-1/4$ and $\psi(2)=1-\gamma>1/4$.
Moreover $1/4\le\psi'(x)\le2$ on $[1,2]$.

Here is an explicit construction of the zero, avoiding a decision whether
an arbitrary computable real is positive.  Maintain rational endpoints
$l,u$ with $\psi(l)\le0\le\psi(u)$ and let $w=u-l$, $m=(l+u)/2$.
Compute a rational $q$ with $|q-\psi(m)|\le\varepsilon=w/64$.
If $q>\varepsilon$, replace $u$ by $m$; if $q<-\varepsilon$, replace $l$
by $m$.  In the remaining, decidable rational case, $|\psi(m)|\le2\varepsilon$.
The derivative lower bound shows that the interval
\[
 [l,u]\cap[m-8\varepsilon,m+8\varepsilon]
\]
still has the required endpoint signs.  Replace $[l,u]$ by that interval.
The width decreases by at least a factor of two each time.  All tests
are rational, all approximations have predetermined error budgets, and
iteration is primitive recursion.  The nested midpoints consequently give
a primitive recursive name $\alpha$ with $1<\alpha<2$ and $\psi(\alpha)=0$.
For the last equality, the bound $\psi'\le2$ bounds the value error by
twice the interval width.  Strict monotonicity of $\psi$ proves uniqueness.
Since $\Gamma'=\Gamma\psi$ and gamma is positive, this point is its unique
global minimum on the positive half-line.  To make this last conclusion
without real-order trichotomy, use the integrated tangent inequality
\[
 g(y)-g(v)-g'(v)(y-v)
 =(y-v)^2\int_0^1(1-t)g''(v+t(y-v))\,dt
 \ge\frac{(y-v)^2}{2b^2}
\]
on any positive bounded interval containing $y$ and $v$.  Substitute
$v=\alpha$ and $g'(\alpha)=0$.  The explicit square gives strictness when
$|y-\alpha|$ has a positive rational lower bound.

The same construction works on every interval $(-n-1,-n)$, $n\ge0$.
For $x$ in that interval, recurrence gives
\begin{equation}\label{eq:negative-digamma}
 \psi(x)=\psi(x+n+2)-\sum_{k=0}^{n+1}\frac1{x+k},\qquad
 \psi'(x)=\psi'(x+n+2)+\sum_{k=0}^{n+1}\frac1{(x+k)^2}>\frac14.
\end{equation}
Take $\delta=1/(4(n+2))$, $l=-n-1+\delta$, and $u=-n-\delta$.
At $l$ the singular term is $-1/\delta$ and at $u$ it is $+1/\delta$.
All other reciprocal terms have absolute value at most $2$; the remaining
digamma term has absolute value at most $1$.  Therefore
\[
 \psi(l)\le-1/\delta+(2n+3)<0,
 \qquad\psi(u)\ge1/\delta-(2n+3)>0.
\]
On $[l,u]$, an upper derivative bound is $2+(n+2)\delta^{-2}$.
The preceding rational algorithm constructs the unique zero $\alpha_n$.
The sign of gamma on the whole interval is $(-1)^{n+1}$, by its recurrence
to the positive interval $(1,2)$.  Thus $\alpha_n$ is a local maximum on
each negative-sign interval and a local minimum on each positive-sign
interval.  In fact these are extrema on the entire indicated intervals:
apply the same tangent inequality to $\log|\Gamma|$, whose second
derivative is $\psi'>0$, and then use the fixed sign of gamma.  This
again avoids splitting an arbitrary real according to its order relative
to the critical point.  No other real critical points exist away from
the poles.

\section{Stirling's formula from finite differences}
\label{sec:stirling-differences}

Define
\[
 A(x)=\log\Gamma(x)-(x-\tfrac12)\log x+x,\qquad
 C=\tfrac12\log(2\pi).
\]
The recurrence for gamma gives the exact finite-difference identity
\[
 d(x):=A(x)-A(x+1)
 =(x+\tfrac12)\log(1+1/x)-1.
\]
With $z=(2x+1)^{-1}\in(0,1)$, the finite geometric expansion for the
logarithm gives
\begin{equation}\label{eq:stirling-d-series}
 d(x)=\sum_{j=1}^{\infty}\frac{z^{2j}}{2j+1}.
\end{equation}
After $j=K$, the error is at most
$z^{2K+2}/((2K+3)(1-z^2))$.  In particular, retaining the improvement in
the coefficient at $j=2$,
\begin{equation}\label{eq:stirling-d-telescoper}
 \frac1{3(2x+1)^2}\le d(x)
 \le\frac1{12x(x+1)}-\frac2{15(2x+1)^4}.
\end{equation}
The upper main term is the difference $1/(12x)-1/(12(x+1))$.
Hence for integers $M>N\ge1$,
\[
 0\le A(N)-A(M)\le\frac1{12N}-\frac1{12M}.
\]
This is a finite sum of inequalities.  It constructs, for example by
sampling $N=2^{p+4}$, a regular primitive recursive name $c$ with
$|A(N)-c|\le1/(12N)$.  No arbitrary monotone-convergence theorem is used.

From \eqref{eq:digamma-elementary-bounds},
\[
 -\frac1{2x}\le A'(x)=\psi(x)-\log x+\frac1{2x}\le0.
\]
Thus for $0\le v\le b$, $|A(N+v)-A(N)|\le b/(2N)$.  In particular the
same $c$ is approached along every fixed translated sequence $N+x$,
with an explicit rational rate.

The duplication formula, proved independently from the finite gamma
products and reflection, gives
\begin{equation}\label{eq:constant-by-duplication}
 A(x)+A(x+\tfrac12)-A(2x)
 =C+\tfrac12-x\log(1+1/(2x)).
\end{equation}
The last two terms lie between $0$ and $1/(8x)$, by
$u-u^2/2\le\log(1+u)\le u$.  Setting $x=N$ and using the preceding
Cauchy estimates proves $|c-C|\le7/(12N)$ for every $N$.  Substituting a
sufficiently large explicit power of two for $N$ proves equality of the
two names.  This determines the constant without choosing a subsequence
or an unknown limiting value.

We now have
\begin{equation}\label{eq:stirling-positive-remainder}
 R(x):=A(x)-C=\sum_{k=0}^{\infty}d(x+k),\qquad
 0\le\sum_{k=N}^{\infty}d(x+k)\le\frac1{12(x+N)}.
\end{equation}
The equality means that for a requested precision the two finite
expressions differ by the displayed tail.  Keeping the first positive
margin in \eqref{eq:stirling-d-telescoper} proves
\begin{equation}\label{eq:stirling-first-bound}
 \frac1{3(2x+1)^2}\le R(x)
 \le\frac1{12x}-\frac2{15(2x+1)^4}<\frac1{12x}
 \qquad(x>0).
\end{equation}
In particular
\begin{equation}\label{eq:stirling-multiplicative-bound}
 1<\frac{\Gamma(x)}{\sqrt{2\pi}\,x^{x-1/2}e^{-x}}<e^{1/(12x)}.
\end{equation}
For $x\ge1$, the relative error is at most $1/(6x)$, since
$e^u-1\le2u$ when $0\le u\le1$.

The sharper Robbins lower bound also follows by rational telescoping.
For $x\ge1$ put
\[
 q(x)=\frac1{12x+1},\qquad
 \lambda(x)=\frac{24x-23}
 {3(2x+1)^2(12x+1)(12x+13)}>0.
\]
The polynomial identity
$(12x+1)(12x+13)-36(2x+1)^2=24x-23$ gives
$d(x)\ge q(x)-q(x+1)+\lambda(x)$.  Summing finitely and using
\eqref{eq:stirling-positive-remainder} proves
\begin{equation}\label{eq:robbins-real}
 \frac1{12x+1}+\lambda(x)\le R(x)
 \le\frac1{12x}-\frac2{15(2x+1)^4}\qquad(x\ge1).
\end{equation}
Taking $x=n$ and using $\Gamma(n+1)=n!$ gives the usual strict factorial
Robbins inequalities.  This argument is the real-variable extension of
the checked finite telescoper in the repository's Stirling draft.

\section{All orders of Stirling's expansion with the sharp signed bound}
\label{sec:all-order-stirling}

The order $m\ge1$ may be fixed externally; the following is then a family
of arithmetic theorems with primitive recursive precision parameters.
Bernoulli polynomials have rational coefficients and are defined without
any analytic generating function: $B_0(t)=1$, $B_r'=rB_{r-1}$, and
$\int_0^1 B_r(t)\,dt=0$ for $r\ge1$.  Antidifferentiating a polynomial
and subtracting its mean is a finite rational algorithm.  Write
$b_r=B_r(0)$.  Uniqueness of these elementary polynomial conditions gives
\[
 B_r(1-t)=(-1)^r B_r(t),\qquad
 B_r(1)=B_r(0)\quad(r\ge2),\qquad b_{2j+1}=0\quad(j\ge1).
\]
For example $B_1=t-1/2$, $B_2=t^2-t+1/6$,
$B_4=t^4-2t^3+t^2-1/30$, and $b_6=1/42$.

\subsection{Finite Euler--Maclaurin, derived rather than assumed}

On $[k,k+1]$ set $\widetilde B_r(t)=B_r(t-k)$, using the common endpoint
value for even $r$.  Integration by parts on that interval gives
\[
 \int_k^{k+1}\widetilde B_1(t)f'(t)\,dt
 =\frac{f(k)+f(k+1)}2-\int_k^{k+1}f(t)\,dt.
\]
Sum over $k=0,\ldots,N-1$ and integrate the last term successively using
$B_r'=rB_{r-1}$.  Odd endpoint contributions vanish and even endpoint
contributions telescope.  One obtains the finite identity
\begin{align}
 \sum_{k=0}^{N}f(k)
 ={}&\int_0^N f(t)\,dt+\frac{f(0)+f(N)}2\notag\\
 &+\sum_{r=1}^{m}\frac{b_{2r}}{(2r)!}
       \bigl(f^{(2r-1)}(N)-f^{(2r-1)}(0)\bigr)
 -\frac1{(2m)!}\int_0^N\widetilde B_{2m}(t)f^{(2m)}(t)\,dt.
 \label{eq:finite-em-gamma}
\end{align}
For the functions used here, integration by parts itself follows by
finite telescoping product increments; derivatives and continuity bounds
are explicit rational functions on a compact interval.  Thus this
identity introduces no nonconstructive calculus principle.

Apply \eqref{eq:finite-em-gamma} to $f(t)=(x+t)^{-1}$ and use the finite
Euler expression for $\psi$.  Since
$f^{(j)}(t)=(-1)^j j!(x+t)^{-j-1}$, all terminal errors vanish with
explicit bounds.  The resulting identity is
\begin{equation}\label{eq:psi-em-exact}
 \psi(x)=\log x-\frac1{2x}
 -\sum_{r=1}^{m}\frac{b_{2r}}{2r x^{2r}}
 +\int_0^\infty\frac{\widetilde B_{2m}(t)}{(x+t)^{2m+1}}\,dt.
\end{equation}
If $M_m$ is the sum of the absolute values of the coefficients of
$B_{2m}$, the last integral after a cutoff $T$ has tail at most
$M_m/(2m(x+T)^{2m})$.  All terminal logarithms are bounded using
$\log(1+x/N)\le b/N$.

Integrate \eqref{eq:psi-em-exact} from $x$ to a finite $Y$, use
$A(Y)-C\to0$ from the preceding explicit Stirling bound, and regroup the
finite double sums.  Integrating the rational kernel gives
\begin{equation}\label{eq:loggamma-em-exact}
 R(x)=\sum_{r=1}^{m}\frac{b_{2r}}{2r(2r-1)x^{2r-1}}
 -\frac1{2m}\int_0^\infty
             \frac{\widetilde B_{2m}(t)}{(x+t)^{2m}}\,dt.
\end{equation}
The omitted rectangular tails in this double-sum argument are at most
$M_m/[2m(2m-1)(x+T)^{2m-1}]$ and
$M_m/[2m(2m-1)Y^{2m-1}]$, respectively.  Thus a general Fubini theorem
is unnecessary.

\subsection{The Bernoulli sign and a finite reverse-Chebyshev argument}

We prove the elementary shape fact needed for the sharp remainder.
On $0<t<1/2$, $B_1(t)<0$.  For $j\ge1$ the polynomial $B_{2j+1}$ vanishes
at $0$ and $1/2$, and
$B_{2j+1}''=(2j+1)(2j)B_{2j-1}$.  Induction therefore proves
\[
 (-1)^jB_{2j-1}(t)>0\qquad(0<t<1/2).
\]
Indeed a function with positive second derivative and zero endpoints is
negative between them, and with negative second derivative it is
positive; for these polynomials this statement follows from the finite
secant inequality proved by integrating derivative increments.

It follows that
\[
 H_m(t)=(-1)^m(B_{2m}(t)-b_{2m})
\]
is symmetric about $1/2$, is zero at $0$ and $1$, is positive in between,
and is strictly increasing on $[0,1/2]$.  Its mean is
\[
 \int_0^1H_m(t)\,dt=(-1)^{m+1}b_{2m}=|b_{2m}|>0.
\]
This also proves the alternating signs of the nonzero even Bernoulli
numbers.

Fix $A>0$ and $q>0$, and put $w(t)=(A+t)^{-q}$.  The paired weight
$W(t)=w(t)+w(1-t)$ is decreasing for $0\le t\le1/2$, since
\[
 W'(t)=-q(A+t)^{-q-1}+q(A+1-t)^{-q-1}\le0.
\]
For an increasing function $H$ and decreasing function $W$ on an interval
of length $L$, the finite double-sum identity for
$(H(u)-H(v))(W(u)-W(v))\le0$ gives
\[
 \int HW\le L^{-1}\left(\int H\right)\left(\int W\right).
\]
Passing from the finite sums uses only their uniform moduli.  Apply this
on $[0,1/2]$ to $H_m$ and the paired weight.  Symmetry yields
\begin{equation}\label{eq:bernoulli-weight-bound}
 0<\int_0^1H_m(t)(A+t)^{-q}\,dt
 <|b_{2m}|\int_0^1(A+t)^{-q}\,dt.
\end{equation}
Strictness has constructive content: $H_m$ is strictly increasing and
$W$ strictly decreasing on the interior, so two disjoint rational
subintervals give an explicit negative margin in the double-sum
identity.  Their rational centers can be fixed, for example at $1/8$
and $3/8$, and continuity bounds shrink the intervals to preserve the
strict gap.  The coefficients of $H_m$ and bounds for $A$ compute the
required widths; no unbounded search is involved.
For an explicit weight gap, if $A\le B$ and $t\in[1/8,3/8]$, then
\[
 -W'(t)\ge\frac{q(q+1)}{4(B+1)^{q+2}},\qquad
 W(1/8)-W(3/8)\ge\frac{q(q+1)}{16(B+1)^{q+2}}.
\]
The other gap, $H_m(3/8)-H_m(1/8)>0$, is an exactly computed rational.

Subtract the first $m-1$ correction terms in
\eqref{eq:loggamma-em-exact}.  Since
$\int_0^\infty(x+t)^{-2m}\,dt=x^{1-2m}/(2m-1)$, the result is
\begin{equation}\label{eq:stirling-remainder-positive-integral}
 E_m(x):=R(x)-\sum_{r=1}^{m-1}
            \frac{b_{2r}}{2r(2r-1)x^{2r-1}}
 =\frac{(-1)^{m+1}}{2m}
       \int_0^\infty\frac{H_m(\{t\})}{(x+t)^{2m}}\,dt.
\end{equation}
Here $H_m(\{t\})$ is implemented piecewise on explicitly indexed unit
intervals; computing the floor of an arbitrary real is not required.
Apply \eqref{eq:bernoulli-weight-bound} on each unit interval and add the
finite inequalities.  The integral tail is explicit.  Keeping the
positive gap from the first interval proves the sharp signed bound
\begin{equation}\label{eq:sharp-stirling-remainder}
 \boxed{\quad
 0<(-1)^{m+1}E_m(x)
 <\frac{|b_{2m}|}{2m(2m-1)x^{2m-1}}
 \quad(x>0).\quad}
\end{equation}
Consequently
\[
 \log\Gamma(x)\sim
 (x-\tfrac12)\log x-x+\tfrac12\log(2\pi)
 +\frac1{12x}-\frac1{360x^3}+\frac1{1260x^5}-\cdots,
\]
where the meaning of every truncation is precisely
\eqref{eq:sharp-stirling-remainder}.  For example,
\[
 \frac1{12x}-\frac1{360x^3}<R(x)<\frac1{12x},\qquad
 0<R(x)-\frac1{12x}+\frac1{360x^3}<\frac1{1260x^5}.
\]

The same argument applied to \eqref{eq:psi-em-exact} gives, if
\[
 D_m(x)=\psi(x)-\log x+\frac1{2x}
             +\sum_{r=1}^{m-1}\frac{b_{2r}}{2r x^{2r}},
\]
the derivative expansion with the sharp signed bound
\begin{equation}\label{eq:sharp-digamma-remainder}
 0<(-1)^mD_m(x)<\frac{|b_{2m}|}{2m x^{2m}}.
\end{equation}

Exponentiating finite truncations gives the multiplicative Stirling
expansion.  Its rational coefficients are determined primitively by
$c_0=1$ and
\[
 n c_n=\sum_{j=1}^{n}j\ell_jc_{n-j},\qquad
 \ell_{2r-1}=\frac{b_{2r}}{2r(2r-1)},\quad \ell_{2r}=0.
\]
Thus the first terms are
\[
 \frac{\Gamma(x)}{\sqrt{2\pi}\,x^{x-1/2}e^{-x}}
 \sim1+\frac1{12x}+\frac1{288x^2}
          -\frac{139}{51840x^3}-\frac{571}{2488320x^4}+\cdots.
\]
For any requested fixed truncation order, choose a log-gamma truncation
with error of higher degree, substitute its finite polynomial into the
finite exponential sum, and bound the two remainders with
$|e^u-e^v|\le e^{\max(|u|,|v|)}|u-v|$ and the geometric majorant for the
exponential Taylor tail.  This is an explicit finite procedure for the
error constant, not an appeal to closure properties of an undefined
asymptotic notation.

\subsection{Primitive recursive quadrature budgets}

Here are sufficient budgets for every improper integral introduced in
this section.  If $P(t)$ is one of the fixed periodic Bernoulli
polynomials, let $M$ be the sum of its absolute coefficient values and
$L$ the sum of the absolute coefficient values of its derivative.  For
$q\ge2$ and $x\ge a>0$, the function
$P(\{t\})(x+t)^{-q}$ is globally Lipschitz on $[0,T]$ with constant
\[
 K=L a^{-q}+qM a^{-q-1}.
\]
The pieces join continuously because the endpoint values of the even
Bernoulli polynomials agree.  A left sum with $J$ equal subintervals has
error at most $KT^2/(2J)$.  Its tail after $T\ge1$ is at most
$M/((q-1)(x+T)^{q-1})\le M/T$.
For target error $2^{-p}$ one may therefore take
\[
 T=\max(1,\lceil 2^{p+3}M\rceil),\qquad
 J=\max(1,\lceil 2^{p+3}KT^2\rceil),
\]
and allocate the remaining error to approximating $x$ in the summands.
The floor symbols in the notation are computed only for rational mesh
points, by integer division.  All loops are bounded, all coefficients
are rational, and all required sums and products are primitive recursive.
For each fixed derivative or asymptotic order, these give the explicit
finite rational identities and inequalities to be formalized inside iPRA.

% Sources for root's bibliography:
% NIST DLMF §§5.5(iv), 5.6(i), 5.11(i)-(iii).
% Repository commit 94a4713736d281e1bfcd0b54e0acb6410136971c:
% bohr_mollerup blob 55a972a2cd46f9b4a5edd70b79598cca3413ccd7;
% stirling blob 6c6f8e17747d567b426b19f0cb2678cba303457f.

\section{A canonical log-gamma and both Binet formulas}
\label{sec:binet}

\subsection{The right-half-plane logarithm}

For $z=x+iy$ with $x\ge a>0$, define
\[
 \operatorname{Log}z
 =\tfrac12\log(x^2+y^2)
    +i\int_0^{y/x}\frac{du}{1+u^2}.
\]
This is a primitive recursive name on each certified bounded rectangle.
Finite differentiation gives its complex derivative $1/z$.
For fixed $x>0$, differentiate
$\exp(\operatorname{Log}(x+iy))/(x+iy)$ in the real variable $y$:
the derivative is zero and its value at $y=0$ is one. The finite
fundamental theorem therefore gives $\exp(\operatorname{Log}z)=z$.
Positive scaling gives
$\operatorname{Log}(cz)=\log c+\operatorname{Log}z$ for $c>0$.
Define
\[
 \mathcal L_N(z)=\log(N!)+z\log N-
                        \sum_{k=0}^N\operatorname{Log}(z+k).
\]
For $N\ge2R$, where $|z|\le R$, its tail is bounded by
$(R^2+R)/N$: rewrite it as
\[
 -\operatorname{Log}z-\gamma_Nz
 -\sum_{k=1}^N\bigl(\operatorname{Log}(1+z/k)-z/k\bigr)
\]
and apply (C2)--(C6). On $|w|\le1/2$ the equality
$\operatorname{Log}(1+w)=w+L(w)$ is explicit: both have derivative
$1/(1+w)$ and value zero at $w=0$, so integrate their difference
along the segment $u\mapsto uw$. Denote the resulting effective limit by
$\mathcal L(z)$. Finite exponentiation and differentiation show
\begin{equation}
 e^{\mathcal L(z)}=\Gamma(z),\qquad
 \mathcal L'(z)=\psi(z),\qquad
 \mathcal L(z+1)-\mathcal L(z)=\operatorname{Log}z.
 \label{eq:canonical-loggamma}
\end{equation}
For the last equality the finite difference is
$\operatorname{Log}z+\log N-\operatorname{Log}(z+N+1)$,
whose last two terms tend explicitly to zero. On the positive axis
$\mathcal L(x)=\log\Gamma(x)$. In general $\mathcal L(z)$ need not be
the principal logarithm of the value $\Gamma(z)$; the normalization is
part of the definition.

\subsection{A positive kernel from reflection}

The differentiated reflection formula at $1/2$ and (C33) give
$3\zeta(2)=\pi^2/2$, hence $\zeta(2)=\pi^2/6$.
Subtract the digamma series (C27) at $z$ and $-z$, and use its recurrence
and reflection formula. Finite reindexing gives
\[
 \pi\cot(\pi z)=\frac1z-
                     2z\sum_{n=1}^{\infty}\frac1{n^2-z^2}.
\]
All derivative and series tails have already been bounded. Substituting
$z=it/(2\pi)$, with $t>0$, proves
\begin{equation}
 k(t):=\frac1t\left(\frac12-\frac1t+\frac1{e^t-1}\right)
       =2\sum_{n=1}^{\infty}\frac1{t^2+4\pi^2n^2}.
 \label{eq:binet-kernel}
\end{equation}
Use the series on the right to define $k(0)=1/12$ as well. It gives
\begin{equation}
 0<k(t)\le\frac1{12},\qquad
 \left|k(t)-2\sum_{n=1}^{N}\frac1{t^2+4\pi^2n^2}\right|
 \le\frac1{2\pi^2N}\le\frac1{8N}\quad(t\ge0).
 \label{eq:binet-kernel-tail}
\end{equation}
For strict positivity the first summand already supplies a rational
lower bound on every bounded interval. Thus the apparent cancellation
at $t=0$ is handled by a regular series rather than subtraction of
singular expressions.

Define
\[
 J(z)=\int_0^\infty e^{-zt}k(t)\dd t\qquad(\Rea z\ge a>0).
\]
The tail after $T$ is at most $e^{-aT}/(12a)$. On $[0,T]$ approximate
the kernel by $N$ summands; its integral error is at most $T/(8N)$.
The derivative of this finite integrand has a rational bound in terms
of $a,R,T$, and therefore a prescribed rational mesh completes the
algorithm. A sufficient coarse choice is a mesh Lipschitz bound
$R+T+1$, since the differentiated kernel series is bounded there by
$T$. Differentiation in $z$ adds powers of $t$; their tails are covered
by the monomial estimates used for Euler's integral.

Put $C=\tfrac12\log(2\pi)$ and
\[
 \mathcal R(z)=\mathcal L(z)-(z-\tfrac12)\operatorname{Log}z+z-C.
\]
Using \eqref{eq:binet-kernel} after multiplying by $t^2$, the finite
geometric sum for $(e^t-1)^{-1}$, and
$\int_0^\infty te^{-(z+n)t}\dd t=(z+n)^{-2}$, we get
\begin{align*}
 J''(z)
 &=\frac1{2z^2}-\frac1z+\sum_{n=1}^{\infty}\frac1{(z+n)^2}\\
 &=\psi'(z)-\frac1z-\frac1{2z^2}=\mathcal R''(z).
\end{align*}
The series tail is at most $1/(a+N)$, so the equality is justified
before passing to the infinite notation.

It remains to fix the possible affine difference; merely comparing
second derivatives would be insufficient. Rectangle-sum comparison for
$t\mapsto(w+t)^{-1}$ in the right half-plane gives
\begin{equation}
 |\psi(w)-\operatorname{Log}w|\le\frac1{\Rea w}.
 \label{eq:complex-psi-bound}
\end{equation}
Indeed the sum of the absolute derivative errors over $[0,N]$ is at
most $1/\Rea w-1/(\Rea w+N)$, the final endpoint contributes at most
$1/(\Rea w+N)$, and
$\log N-\operatorname{Log}(w+N)$ has an explicit vanishing bound.
Consequently $|\mathcal R'(w)|\le3/(2\Rea w)$. Integrate vertically
from $N+\Rea z$ to $N+z$ and use the proved real Stirling estimate:
\[
 |\mathcal R(N+z)|\le
 \frac{1/12+(3/2)|\Ima z|}{N+\Rea z},\qquad
 |J(N+z)|\le\frac1{12(N+\Rea z)}.
\]
The function $D=\mathcal R-J$ has zero second derivative, so finite
fundamental-theorem integration on a horizontal segment yields
\[
 D(z)=2D(z+N)-D(z+2N).
\]
The displayed bounds make the right side arbitrarily small with a
primitive recursive choice of $N$. Hence $D(z)=0$.

\begin{theorem}[Binet's first formula]
On every certified bounded region in $\Rea z>0$,
\[
 \boxed{\displaystyle
 \mathcal L(z)=(z-\tfrac12)\operatorname{Log}z-z
 +\tfrac12\log(2\pi)
 +\int_0^\infty
 \left(\frac12-\frac1t+\frac1{e^t-1}\right)
                 \frac{e^{-zt}}t\dd t.}
\]
At zero the integrand is interpreted using the regular kernel
\eqref{eq:binet-kernel}; all convergence estimates above are uniform
on the certified region.
\end{theorem}

\subsection{Binet's second formula without a contour deformation}

For real $t\ge0$ and $\Rea z>0$, specify the branch
\[
 \operatorname{atan}(t/z):=
                 \int_0^t\frac{z}{z^2+u^2}\dd u.
\]
The denominator factors $z+iu$ and $z-iu$ are each apart from zero
by $\Rea z$. The elementary Laplace integral for cosine gives
\[
 \int_0^\infty e^{-zs}\cos(ts)\dd s=\frac{z}{z^2+t^2}.
\]
It is verified by writing cosine as two exponentials and integrating
their finite antiderivatives. The sine integral on the right of the
next display has $t$-derivative equal to this cosine integral: on a
finite interval this is finite differentiation, and its derivative
tail is at most $e^{-aS}/a$. Comparing its value at $t=0$ with the
defining arctangent integral proves
\begin{equation}
 \operatorname{atan}(t/z)=
       \int_0^\infty e^{-zs}\frac{\sin(ts)}s\dd s,
 \qquad |\operatorname{atan}(t/z)|\le t/a
 \quad(\Rea z\ge a).
 \label{eq:atan-laplace}
\end{equation}
At $s=0$ use $\sin(ts)/s=t\operatorname{sinc}(ts)$. The estimate
$|\sin(ts)|\le ts$ for nonnegative real $s,t$ supplies both the
integrable majorant and the derivative justification.

Insert \eqref{eq:atan-laplace} into
$2\int_0^\infty\operatorname{atan}(t/z)/(e^{2\pi t}-1)\dd t$.
The resulting double integral has absolute majorant
\[
 \frac{2e^{-as}t}{e^{2\pi t}-1}.
\]
This is an explicit finite-sum argument: the $s>S$ tail is at most
$4e^{-aS}/a$, and the $t>T\ge1$ tail is at most
$4(T+1)e^{-T}/a$. The small strips $s<\eta$ and $t<\eta$
contribute at most $4\eta$ and $\eta/(\pi a)$ respectively.
These coarse bounds follow by splitting the remaining integral at $1$,
using $e^{2\pi t}-1\ge2\pi t$ below $1$ and
$(e^{2\pi t}-1)^{-1}\le2e^{-t}$ above it.
On the remaining finite rectangle the order of summation is a finite
rearrangement.

Expand the geometric denominator in the inner $t$ integral. Its $n$th
Laplace sine integral, after the regularized division by $s$, is
\[
 \int_0^\infty e^{-2\pi nt}\frac{\sin(st)}s\dd t
              =\frac1{s^2+4\pi^2n^2}.
\]
The tail after $n=N$ is at most $1/(4\pi^2N)$, uniformly in real
$s\ge0$, since $|\sin(st)/s|\le t$. Thus the inner sum is precisely
$k(s)/2$, proving
\begin{theorem}[Binet's second formula]
With the preceding right-half-plane branches,
\[
 \boxed{\displaystyle
 \mathcal L(z)=(z-\tfrac12)\operatorname{Log}z-z
 +\tfrac12\log(2\pi)
 +2\int_0^\infty
       \frac{\operatorname{atan}(t/z)}{e^{2\pi t}-1}\dd t.}
\]
All exchanges have been reduced to finite rectangles with the stated
tails. In particular no contour integral or branch continuation theorem
is needed.
\end{theorem}

For real $x>0$, the positive-kernel formula immediately proves the
additional consequence
\[
 (-1)^j R^{(j)}(x)=\int_0^\infty t^je^{-xt}k(t)\dd t>0
 \qquad(j\in\N).
\]
A positive rational margin follows by restricting the integral to
$[1,2]$ and using its first positive summand. This is complete
monotonicity of the Stirling remainder with an explicit constructive
meaning.

\section{Log-gamma integrals, Mellin transforms, and applications}
\label{sec:further}

\subsection{Raabe's formula and the unit-interval integral}

Let $J_s=\int_0^{\pi/2}\log(\sin t)\dd t$. These endpoint integrals
have effective tails: by $2t/\pi\le\sin t\le t$, an omitted interval
of length $\eta$ contributes at most
$\eta(1+|\log\eta|+\log2)$ in absolute value. Reflection about
$\pi/2$ gives the same bound for cosine. Substitution and the double
angle identity therefore give, first on truncated intervals,
\[
 2J_s=\int_0^{\pi/2}\log(\sin t\cos t)\dd t
     =-\tfrac\pi2\log2+\tfrac12\int_0^\pi\log(\sin u)\dd u
     =-\tfrac\pi2\log2+J_s.
\]
Hence $\int_0^1\log(\sin\pi t)\dd t=-\log2$.
The endpoint behavior
$\log\Gamma(t)=-\log t+\log\Gamma(1+t)$ is equally effective:
for $0<t\le1/2$, (C37)--(C38) give
$|\log\Gamma(1+t)|\le2$, so the omitted gamma-log integral is at
most $\eta(3+|\log\eta|)$. With $\eta=2^{-h}$, the choice
$h=4(p+10)$ makes each such tail less than $2^{-p}$.

Take real logarithms in the reflection identity and integrate on
$[\eta,1-\eta]$. Positivity licenses the logarithms; the preceding
tails allow the endpoints to be restored. We obtain
\begin{equation}
 \boxed{\displaystyle
 \int_0^1\log\Gamma(t)\dd t=\tfrac12\log(2\pi).}
 \label{eq:unit-loggamma}
\end{equation}
For $x>0$ put $I(x)=\int_x^{x+1}\log\Gamma(t)\dd t$.
The finite fundamental theorem and the recurrence give
$I'(x)=\log\Gamma(x+1)-\log\Gamma(x)=\log x$.
Thus $I(x)-x\log x+x$ is constant. Letting $x$ tend to zero with the
explicit endpoint estimates just proved identifies the constant in
\eqref{eq:unit-loggamma}. This proves Raabe's formula
\begin{equation}
 \boxed{\displaystyle
 \int_x^{x+1}\log\Gamma(t)\dd t
 =x\log x-x+\tfrac12\log(2\pi)\qquad(x>0).}
 \label{eq:raabe}
\end{equation}
The recurrence also gives
$\int_0^1\log\Gamma(1+t)\dd t=\tfrac12\log(2\pi)-1$.

\subsection{The gamma--zeta Mellin identity}

For $\Rea s\ge1+1/A$, define $\zeta(s)$ by its Dirichlet sums.
Their absolute tail is at most $A N^{-1/A}$, by integrating
$t^{-1-1/A}$ from $N$ onwards. In particular a cutoff
$N\ge(2^{p+4}A)^A$ gives a primitive recursive name. The integral
\[
 \int_0^\infty\frac{t^{s-1}}{e^t-1}\dd t
\]
has small-endpoint bound $A\eta^{1/A}$, since $e^t-1\ge t$,
and large-endpoint bound twice the gamma tail, since
$(e^t-1)^{-1}\le2e^{-t}$ for $t\ge1$.
For a finite $N$, integrate the exact finite geometric sum
$\sum_{n=1}^N e^{-nt}$ against $t^{s-1}$ and use the positive real
scaling substitution. The result is
$\Gamma(s)\sum_{n=1}^N n^{-s}$. The omitted integral is absolutely
bounded by
\[
 \Gamma(\Rea s)\sum_{n>N}n^{-\Rea s}
 \le (A+2^{R+1}R!)A N^{-1/A},\quad |s|\le R.
\]
This last comparison is established first for a finite sum of the
positive majorants and then by their stated geometric and integral
tails; it is not an appeal to a measure-theoretic interchange theorem.
It proves
\begin{equation}
 \boxed{\displaystyle
 \Gamma(s)\zeta(s)=\int_0^\infty\frac{t^{s-1}}{e^t-1}\dd t,
 \qquad\Rea s>1.}
 \label{eq:gamma-zeta-mellin}
\end{equation}
For $s=2$ this includes
$\int_0^\infty t/(e^t-1)\dd t=\pi^2/6$.
The zeta functional equation is a separate theorem and is not being
inferred from this half-plane identity.

\subsection{Wallis, Gaussian moments, and factorial interpolation}

At $z=1/2$ the sine product gives
\[
 \frac\pi2=\prod_{k=1}^{\infty}\frac{4k^2}{4k^2-1}.
\]
If $W_N$ is its finite product, then
$0\le\pi/2-W_N\le2/N$. To see the rate, the remaining logarithm is
at most $\sum_{k>N}2/(4k^2)\le1/(2N)$, and
$e^u-1\le2u$ for $0\le u\le1$. Use $W_N<\pi/2<2$.

The quadratic integral substitution and the half-integer values give,
for $c>0$ and $n\in\N$,
\[
 \int_{-\infty}^{\infty}t^{2n}e^{-ct^2}\dd t
   =\frac{(2n)!\sqrt\pi}{4^n n!c^{n+1/2}},\qquad
 \int_{-\infty}^{\infty}t^{2n+1}e^{-ct^2}\dd t=0.
\]
The odd integral is absolutely convergent, and its two finite symmetric
halves cancel exactly. In particular the Gaussian integral is
$\sqrt{\pi/c}$. Likewise the positive function
$p_x(t)=t^{x-1}e^{-t}/\Gamma(x)$ has integral one, and its finite
integer moments are
\[
 \int_0^\infty t^n p_x(t)\dd t=(x)_n.
\]
Thus its mean and variance are $x$ and $x$, by the cases $n=1,2$.
These statements are integral identities with the already proved
cutoffs; no additional probability-theoretic foundation is assumed.

Two equations mentioned in the reference article reduce to algebra.
For $z>1/2$,
$\Gamma(z-1/2)=\Gamma(z+1/2)$ holds exactly when $z=3/2$,
by recurrence and positive cancellation. For positive $z$ with
$z-1$ apart from the poles, the equality
$\Gamma(z-1)=\Gamma(z+1)$ is equivalent to $z(z-1)=1$.
Its unique positive solution is
$\phi=(1+\sqrt5)/2$. In fact the equation itself forces
$z-1=1/z>0$, and then positive square-root uniqueness solves
$(2z-1)^2=5$; there is no real-order trichotomy step.

Finally, normalization and recurrence alone do not characterize gamma.
For a nonzero rational $c$,
\[
 F_c(z)=\Gamma(z)\exp(c\sin(2\pi z))
\]
satisfies the same recurrence and factorial values. It differs from
gamma at $z=1/4$. Its analytic and arithmetic presentation follows from
the functions already constructed. This supplies explicit alternative
interpolations and explains why the convexity hypothesis in the
Bohr--Mollerup theorem carries mathematical content.

\section{Executable arithmetic and a complete Lean formalization route}
\label{sec:formalization}

\subsection{What has actually been checked}

The companion \code{lean/GammaIPRA.lean} was compiled with Lean
4.32.0, using only \code{Std}. It contains exact rational complex
arithmetic, finite elementary approximations, Euler products, numerical
budget functions, representations of names, and analytic certificate
interfaces. Fourteen named elementary arithmetic lemmas are proved.
The saved \code{\#print axioms} audit reports no axiom dependencies for
each of these fourteen lemmas. Two additional closed arithmetic examples
are checked by kernel reduction. The file contains no admitted proofs
or axiom declarations.

\emph{This does not constitute a formal proof of gamma analysis.}
The analytic certificate structures are uninstantiated obligations.
In particular, the file does not yet prove that its elementary
truncation lengths imply the corresponding analytic error bounds, and
it does not yet produce \code{PRAAx} derivations of the displayed
estimates. These two substantial remaining tasks are specified below.
The human-readable proofs in the preceding sections are the mathematical
argument; the Lean files are an executable foundation and a precise
formalization outline.

\subsection{The exact rational demonstration and its bracket theorem}

The program \code{code/gamma\_brackets.py} evaluates gamma on positive
rational arguments without floating point. We include the proof of its
particular enclosure method so its certificates have a mathematical
interpretation independent of a numerical gamma library.

For a rational $0<r\le1$ set
\[
 A_n(r)=\frac{n!n^r}{r(r+1)\cdots(r+n)},\qquad
 B_n(r)=\frac{(n-1)!n^r}{r(r+1)\cdots(r+n-1)}.
\]
Then
\begin{equation}
 A_n(r)\le\Gamma(r)\le B_n(r),\qquad
 0\le B_n(r)-A_n(r)\le1/n.
 \label{eq:computational-strip-brackets}
\end{equation}
Here is a short complete proof, also explaining the bracket construction
in the inspected Bohr--Mollerup draft~\cite{bmrepo}. For $0<r\le1$,
the concavity inequality $t^r\le1+r(t-1)$ for $t>0$ follows by
integrating the decreasing derivative $r t^{r-1}$ from $1$ to $t$.
Equivalently for rational $r$ it can be cleared to positive integer
power comparisons. Exact cancellation gives
\[
 \frac{A_{n+1}}{A_n}=
 \frac{n+1}{n+1+r}(1+1/n)^r,
 \qquad
 \frac{B_{n+1}}{B_n}=
 \frac{n}{n+r}(1+1/n)^r.
\]
Concavity at $t=1+1/n$ makes the second ratio at most one.
Concavity at $t=n/(n+1)$, followed by positive reciprocation, gives
\[
 (1+1/n)^r\ge(1-r/(n+1))^{-1}\ge1+r/(n+1),
\]
so the first ratio is at least one. Since
$B_n=(1+r/n)A_n$, the brackets are nested. Also
$A_n\le B_n\le B_1=1/r$, and hence
$B_n-A_n=(r/n)A_n\le1/n$. Their common effective limit is the
Euler limit already identified with gamma, proving
\eqref{eq:computational-strip-brackets}.

For a positive rational $x$, integer division computes
$x=r+m$ with $0<r\le1$ and $m\in\N$. This uses a floor only for
a rational input, not for a named real. Put $S=(r)_m>0$. The recurrence
gives $\Gamma(x)=S\Gamma(r)$. For target width $2^{-p}$ the program
chooses
\[
 n=2^{p+2}\max(1,\ceil S).
\]
Write $r=u/v$ in lowest terms, and write $A_n=c_A n^r$,
$B_n=c_B n^r$ with rational $c_A,c_B>0$. It computes a dyadic root
enclosure of $n^r$ by bounded binary search for the integer $k$ such that
\begin{equation}
 k^v\le n^u2^{vL}<(k+1)^v.\label{eq:root-certificate}
\end{equation}
An initial upper endpoint is a power of two whose exponent is obtained
by integer bit length; that many bisection iterations plus one suffice.
Choose
\[
 L=p+3+\operatorname{bitlength}
                    \bigl(\max(1,\ceil{Sc_B})\bigr).
\]
Then $Sc_B2^{-L}\le2^{-p-3}$, while $S/n\le2^{-p-2}$.
The interval
\[
 \left[Sc_A\frac{k}{2^L},\quad
       Sc_B\frac{k+1}{2^L}\right]
\]
contains $\Gamma(x)$ and has width less than $2^{-p}$.
Every quantity in \eqref{eq:root-certificate} is an integer and is
stored in the evidence file. The verification routine reconstructs the
rational coefficients, checks the integer inequalities and interval
width, and checks that the displayed decimal endpoints round outwards.
It does not test the answer against a floating-point gamma oracle.

The supplied run uses eight rational arguments and $p=8$; for example
it gives the rigorous enclosures
\[
 \Gamma(1/2)\in[1.7718050888,1.7726837525],\qquad
 \Gamma(3/2)\in[0.8859025444,0.8863418763].
\]
These finite certificates check implementation arithmetic. Their
interpretation as gamma enclosures uses the bracket theorem above;
samples are not proofs of universal analytic identities.

\subsection{A fully specified rational-input gamma evaluator}

For the complex construction there is also a direct, deliberately
inefficient, closed primitive recursive evaluator. This makes the
formalization target concrete even though the economical strip program
is preferable for examples.

Supply $R\ge1$, $D\in\N$, $\norm z_1\le R$, and the rational tests
$|z+j|\ge1/(D+1)$ for $0\le j\le2R$. Set
\[
 M=M_R,\quad B=R^2+R,\quad
 L=(D+1)^{2R+1}(2R)!3^{2R^2+2R};
 \qquad\lambda=1/L.
\]
For requested precision $p$, choose
\begin{align*}
 N&=\max\{2R,B,3MBL^2 2^{p+6}\},\\
 A&=R(1+R)^N,\qquad T=R(N+1)+1,\\
 H&=A\bigl((R+1)3^{T+1}+1\bigr),\qquad
 \ell=p+H+2L+12.
\end{align*}
Compute $q=\log N$ using the atanh sum with
$N(\ell+2N+4)$ terms, so its error is at most $2^{-\ell}$.
Compute $E=\exp(-zq)$ by the Taylor sum through degree
$4T+\ell+6$. Form the rational complex number
\[
 U=z\prod_{j=1}^N(1+z/j)\,E,
\]
and return its rational inverse. Every loop length is now given by a
closed primitive recursive expression.

To verify the precision budget, the exact polynomial head has
$1$-norm at most $A$. The propagated logarithm and exponential errors
are at most $H2^{-\ell}\le2^{-p-12}/L^2$.
The lower bound for $Q_N$ is $1/L$, so $U$ is apart from zero by at
least $1/(2L)$. Inversion amplifies the $1$-norm error by at most
$4L^2$. Formula (C15) gives the Euler tail at most $2^{-p-6}$ in
complex modulus, hence at most $2^{-p-5}$ in the $1$-norm. The total
error is less than $2^{-p-4}$. The spare margin gives the regular-name
inequality and accommodates the fixed shifts in the equality convention.
Changing $R,D$ to any larger valid domain certificate yields an equal
name by comparing the two finite Euler approximants with the same
finite-tail bound. No search for a modulus occurs.

\subsection{Definitions needed by the full Lean development}

The host-level interfaces in \code{GammaIPRA.lean} are convenience
types for organizing the proof. They are not asserted to be internal
sorts of \ipra. The following definitions fix the intended completion.

\begin{longtable}{@{}>{\raggedright\arraybackslash}p{0.24\textwidth}>{\raggedright\arraybackslash}p{0.70\textwidth}@{}}
\toprule
Object & Definition and required data\\\midrule\endhead
Rational code & Signed integer numerator and positive natural denominator;
every relation is cleared to integer multiplication and comparison.\\
Complex rational & A pair of rational codes with the operations in
\code{CQ}; squared modulus is $u^2+v^2$, and the $1$-norm is $|u|+|v|$.\\
Fixed PR implementation & A repository \code{PR k} expression and its
defining-equation proofs for the chosen routine. The associated function
symbol is used directly; there is no universal evaluation symbol.\\
Regular name & One fixed PR rational-output function, with the
$2^{-p}+2^{-q}$ Cauchy certificate. \code{NameEq} is equality at every
precision with the stated two-stage slack.\\
Domain data & Integer radius $R$, positive rational margin, and finite
pole-distance tests. For a named input, refine its approximation until
the supplied margin is preserved at half its size.\\
Gamma evaluator & The closed evaluator just specified, or its equivalent
reciprocal-product name. Domain-independent equality is a theorem.\\
Euler constant and zeta values & Harmonic-minus-log with (C2), and finite
reciprocal-power sums with (C17).\\
Analytic disk certificate & Coefficient names, a strictly larger rational
disk, coefficient majorant, finite polynomial degree schedule, and
coefficient-rounding and approximation bounds. This is
\code{DiskAnalyticCertificate}.\\
Derivative certificate & A quadratic remainder bound at rational inputs,
with adjustable rational evaluation precision, as in
\code{DerivativeRemainderCertificate}; higher orders iterate finite
polynomial coefficient operations.\\
Integral certificate & Explicit endpoint-tail functions, a compact mesh
error bound, and the finite Riemann-sum program. The integral name is
their specified diagonal evaluation.\\
Derived functions & $\psi=\Gamma'/\Gamma$ with a denominator margin;
polygamma by (C29); beta by (B1); log-gamma by
\eqref{eq:canonical-loggamma}; incomplete gamma by its two truncated
integrals; Bernoulli polynomials by finite antidifferentiation and zero
mean. All have definitions before any theorem invokes them.\\
\bottomrule
\end{longtable}

The analytic certificate in the Lean file uses a coefficient bound
$|c_n|\le M r^{-n}$ on an outer radius $r$ larger than the working
radius. This is enough to prove geometric convergence on the smaller
disk. It must not be replaced by the assertion that some arbitrary
sequence of polynomials approximates the function: that weaker
assertion would not establish analyticity.

\subsection{The proof dependency order and concrete remaining obligations}

\begin{longtable}{@{}>{\raggedright\arraybackslash}p{0.25\textwidth}>{\raggedright\arraybackslash}p{0.42\textwidth}>{\raggedright\arraybackslash}p{0.26\textwidth}@{}}
\toprule
Planned module & Lemmas to implement & Finite invariant / induction\\\midrule\endhead
\code{RatCode} & Sign, fraction, gcd and division implementations;
norm inequalities and error propagation. & Cleared integer identities;
bounded remainder invariants.\\
\code{Elementary} & Exponential and logarithm tails; positive roots;
sine, cosine, arctangent and $\pi$. & Finite sum or product plus a
geometric/factorial tail term.\\
\code{Names} & Regularity, equivalence, arithmetic operations,
margin-preserving inversion, local continuity extension. & Fix the
output precision first; all stage comparisons are rational.\\
\code{EulerProducts} & (C1)--(C15); instantiate the reciprocal and
gamma names and \code{EulerErrorCertificate}. & Partial sums of
$k^{-2}$ and telescoping product error.\\
\code{AnalyticGamma} & (C16)--(C26); instantiate local analytic and
derivative certificates; simple zeros and poles. & Finite coefficient
convolution, binomial translation, geometric inversion.\\
\code{EulerIntegral} & Mesh lemmas and
\eqref{eq:integral-product-error}; derivative integrals and substitutions.
& Compact telescoping sums, explicit small and large tail budgets.\\
\code{FunctionalRelations} & Sine factorization, (T4), reflection,
Gauss, duplication and beta formulas. & Polynomial division by known
separated roots, finite denominator regrouping, beta recurrence.\\
\code{DifferentialGamma} & (C27)--(C45), differentiated reflection and
multiplication, rational-argument derivative values. & Finite power
coefficients and rational reciprocal tails.\\
\code{RealGamma} & Strong convexity, Bohr--Mollerup, Wendel/Gautschi,
ratio estimates and all critical-point algorithms. & Secant inequalities
and rational brackets, preserving endpoint signs.\\
\code{Stirling} & Finite differences, normalization by duplication,
Robbins, finite Euler--Maclaurin and signed remainder. & Finite
telescoper; for each fixed order, rational Bernoulli polynomial
identities and finite reverse-Chebyshev sums.\\
\code{GammaIntegrals} & Binet, Raabe, Mellin, Gaussian and incomplete
gamma consequences. & Finite rectangles and geometric kernels,
with every tail allocated before exchange.\\
\code{ObjectCertificates} & Reify the preceding finite statements as
\code{Fm PTm}; produce \code{Prf PRAAx} derivations. & Use
\code{PRA.qfInd} only on the encoded rational invariants.\\
\bottomrule
\end{longtable}

The entries describe actual proof obligations. None is silently supplied
by an axiom with a name such as ``gamma correctness.'' The earlier
sections give the required finite identity, error bound, or polynomial
argument for each entry. A proof in Mathlib's completed real or complex
field could be added for exposition, but cannot replace
\code{ObjectCertificates} when the target is provability in iPRA.

\subsection{One explicit object-theory target}

Let $\mathsf A_{R,D}(z,p)$ denote the rational evaluator above.
An internal functional-equation target is the integer-coded formula
\[
 \begin{aligned}
 &\forall R,D,z,p\;\mathrm{Domain}(R,D,z,z+1)\ \longrightarrow\\
 &\quad\left\|
 \mathsf A_{R+1,D}(z+1,p+R+6)
 -z\,\mathsf A_{R+1,D}(z,p+R+6)
 \right\|_1\le2^{-p}.
 \end{aligned}
\]
Here $z$ abbreviates two rational codes; \code{Domain} is the finite
list of squared-distance and denominator comparisons, expanded by a
primitive recursive Boolean conjunction. The displayed inequality is
also a Boolean comparison after clearing positive denominators. The
proof starts from the exact finite relation
\[
 (z+N+1)\Gamma_N(z+1)=Nz\Gamma_N(z),
\]
proved by the finite rising-product identity, and bounds its residual
using $|z+N+1|\ge N/2$ at $N\ge2(R+1)$. Substituting the explicitly
chosen common index and the arithmetic evaluation errors proves the
integer-coded formula by quantifier-free induction and equality rules.
Comparing arbitrary two cutoff indices gives the Cauchy theorem; it
does not require the theory to quantify over a limit object.

Reflection, for example, has exactly the same form with
$\sin(\pi z)\Gamma(z)\Gamma(1-z)-\pi$ on the left and the finite
sine estimate (T4) in place of the recurrence residual. Clearing the
identity before dividing is preferable in the formal code. The integral
statements compare two explicitly defined rational algorithms with a
$2^{-p}$ error, not a bare Lean \code{HasIntegral} assertion.

Once an ordinary arithmetic sequent $h:\Gamma\vdash[\code{PRAAx}]\varphi$
has been built, the inspected repository theorem gives its native
translation directly:
\[
 \code{pra\_iff\_native.mp}\ h
 :\ \code{NPrf}\bigl(\Gamma.\code{map}\ \code{NFm.ofPRA}\bigr)
                    (\code{NFm.ofPRA}\ \varphi).
\]
This is where the existing ConciseIPRA result is reused. The bridge
translates a supplied proof; it does not automatically prove new gamma
estimates from their truth in the host natural numbers.

\subsection{Induction and uniformity audit}

The rational evaluator, the critical-point algorithms, finite coefficient
operations, and all stated cutoff formulas are primitive recursive.
Proofs of finite algebraic recursions use natural-number parameters
internally. Assertions indexed by derivative or asymptotic order may
always be read as \emph{a family of iPRA proofs, one for each fixed
order}. The present outline does not silently strengthen this to
induction over an unrestricted analytic predicate. A future uniform
object theorem in the order parameter requires reifying the finite
polynomial invariant, exactly as specified in the module table.

No proof uses an unbounded search for a zero, excluded middle for an
arbitrary real comparison, countable choice, an arbitrary Cauchy
completeness axiom, or the floor of an arbitrary real. Unknown signs in
root finding use the quantitative third branch of the bisection
algorithm. Every complex division carries an explicit margin; all
global assertions about gamma at its poles use the analytic reciprocal
or a punctured Laurent expression.

\section{Coverage ledger and the exact boundary of the result}
\label{sec:coverage}

The requested Wikipedia and MathWorld articles~\cite{wiki,mathworld}
contain definitions, analytic identities, numerical tables, historical
remarks, implementation discussions, and deep arithmetic assertions.
The following ledger records \emph{62 result families developed here}.
Some entries contain several related equations; the number is a coverage
index, not a count of independent new theorems. The standard formulas
were cross-checked against the corresponding NIST chapter~\cite{dlmf}.
The contribution here is the finite constructive argument and its
arithmetic interpretation, not novelty of the gamma identities.

All entries have the fixed-name, finite-precision meaning explained in
\cref{sec:foundations}. ``Complex'' includes explicit separation from
every denominator's zero set. ``Positive'' includes a supplied positive
rational lower bound. Individual derivative and asymptotic orders can
be read as externally fixed, as explained in
\cref{sec:formalization}. The human-readable proofs are present; the
full Lean object-theory formalization remains the specified work.

\small
\begin{longtable}{@{}r >{\raggedright\arraybackslash}p{0.52\textwidth}>{\raggedright\arraybackslash}p{0.34\textwidth}@{}}
\toprule
ID & Result family & Proof / domain qualification\\\midrule\endhead
1 & Euler's constant and its harmonic-log error bound & (C1)--(C2), error $1/N$.\\
2 & Euler's limit for gamma & (C15), bounded pole-separated inputs.\\
3 & Weierstrass product for $1/\Gamma$ & (C7)--(C11), every disk.\\
4 & Entire analyticity of the reciprocal & (C16)--(C19), coefficient majorants.\\
5 & Reciprocal zeros exactly $0,-1,-2,\ldots$, all simple & (C22)--(C23) and finite root location.\\
6 & Gamma has no zeros & Positive lower bound $1/M_R$.\\
7 & Meromorphic gamma; simple poles and residues $(-1)^n/n!$ & (C24)--(C25), punctured neighborhoods.\\
8 & Functional equation and normalization & (C20)--(C21).\\
9 & Integer factorial values & (C21), finite recurrence.\\
10 & Complex conjugation and real values on the real domain & (C26).\\
11 & Euler's integral and absolute convergence & \cref{sec:integrals}, explicit endpoint tails.\\
12 & Positivity and computable upper and lower bounds & \eqref{eq:gamma-positive}--\eqref{eq:gamma-real-bound}.\\
13 & $|\Gamma(\sigma+iy)|\le\Gamma(\sigma)$ & Integral triangle inequality, $\sigma>0$.\\
14 & All derivatives as integrals with logarithmic factors & \cref{sec:integrals}, fixed derivative order.\\
15 & Exponential, unit-interval, and quadratic integral forms & \eqref{eq:exponential-substitution}--\eqref{eq:quadratic-substitution}.\\
16 & General positive scaling integral & \eqref{eq:general-scaling}, $c,r>0$.\\
17 & Lower and upper incomplete gamma: sum, recurrence and integer formula & End of \cref{sec:integrals}.\\
18 & Euler reflection and $\Gamma(z)\Gamma(-z)$ & (R1)--(R3), explicit sine margin.\\
19 & Gauss multiplication & (M1)--(M5), arbitrary positive integer multiplier.\\
20 & Legendre duplication & (M6), independent of Stirling.\\
21 & Triplication & (M7).\\
22 & Beta--gamma relation & (B1)--(B8), both real parts positive.\\
23 & Beta symmetry and recurrence & (B2)--(B5).\\
24 & Beta integral evaluations and trigonometric forms & (B9)--(B13), stated half-plane margins.\\
25 & $\Gamma(1/2)=\sqrt\pi$ & (R4), positivity selects the root.\\
26 & All positive and negative half-integer values & (V1).\\
27 & Third, quarter, fifth, sixth and twelfth relations & (V2)--(V6), positive roots.\\
28 & Product of all fractional arguments $r/m$ & (V7).\\
29 & Rising factorial and generalized binomial identities & (V8)--(V9), polynomial cancellation at apparent poles.\\
30 & Imaginary and half-shifted imaginary modulus formulas & (VM1)--(VM5), with a zero-safe version.\\
31 & Digamma partial fractions and their error & (C27)--(C28).\\
32 & All polygamma sums, recurrences, and sign patterns & (C29)--(C35).\\
33 & Integer and half-integer derivative values & (C32)--(C34), $\Gamma'(1)=-\gamma$.\\
34 & Differentiated reflection and multiplication & (D1)--(D3).\\
35 & Digamma values at thirds, quarters and sixths & (D4).\\
36 & Gamma derivatives by finite polynomial recursion & (C36).\\
37 & Taylor expansion at one and local log-gamma series & (C37)--(C41), disk radius $<1$.\\
38 & Entire reciprocal Maclaurin series and coefficient recurrence & (C19), (C42).\\
39 & Laurent constant and quantitative remainder at every pole & (C43).\\
40 & Expansions about every positive integer & (C44)--(C45), disk radius $<n$.\\
41 & Vertical-line modulus infinite product & (C46), real part separated from poles.\\
42 & Balanced finite-parameter infinite products & (C47)--(C48).\\
43 & Strict log-convexity and strict convexity & \eqref{eq:strong-log-convexity}, quantitative gaps.\\
44 & Effective Bohr--Mollerup uniqueness & \eqref{eq:bm-width-general}, fixed function presentations.\\
45 & Wendel and Gautschi inequalities & \eqref{eq:wendel-constructive}--\eqref{eq:gautschi-constructive}.\\
46 & Gamma-ratio asymptotics with a rate & \eqref{eq:general-gamma-ratio}, real shifts.\\
47 & Unique positive minimum and a PR algorithm for its location & \cref{sec:gamma-critical-points}.\\
48 & Sign alternation and one extremum between every two negative poles & \cref{sec:gamma-critical-points}.\\
49 & Stirling leading term and explicit relative error & \eqref{eq:stirling-first-bound}--\eqref{eq:stirling-multiplicative-bound}, $x>0$.\\
50 & Robbins factorial inequalities, and a real-variable strengthening & \eqref{eq:robbins-real}, $x\ge1$.\\
51 & Every order of the log-Stirling expansion with sharp signed remainder & \eqref{eq:sharp-stirling-remainder}, fixed order, real $x>0$.\\
52 & Every order of the digamma asymptotic expansion & \eqref{eq:sharp-digamma-remainder}, fixed order, real $x>0$.\\
53 & Multiplicative Stirling coefficients and an error-producing algorithm & End of \cref{sec:all-order-stirling}.\\
54 & Canonical analytic log-gamma on the right half-plane & \eqref{eq:canonical-loggamma}; branch specified.\\
55 & Binet's first and second formulas & \cref{sec:binet}, complex $\Rea z>0$.\\
56 & Complete monotonicity of the real Stirling remainder & Positive Laplace kernel in \cref{sec:binet}.\\
57 & Raabe and the unit-interval log-gamma integrals & \eqref{eq:unit-loggamma}--\eqref{eq:raabe}.\\
58 & Gamma--zeta Mellin identity & \eqref{eq:gamma-zeta-mellin}, $\Rea s>1$.\\
59 & Wallis product and a rational convergence rate & \cref{sec:further}.\\
60 & Gaussian moments and gamma-density normalization and moments & \cref{sec:further}.\\
61 & The two gamma-crossing equations and the golden ratio solution & \cref{sec:further}, valid gamma domains.\\
62 & Explicit alternative factorial interpolations satisfying recurrence & $\Gamma(z)e^{c\sin(2\pi z)}$; \cref{sec:further}.\\
\bottomrule
\end{longtable}
\normalsize

\subsection{Statements deliberately not claimed}

The scope boundary is mathematical, not a claim that the omitted
statements are impossible in iPRA. We do not provide proofs here of
H\"older's differential-transcendence theorem; irrationality,
transcendence, or algebraic-independence results for special gamma values;
elliptic-integral singular-value evaluations; Kummer's Fourier series;
the displayed continued fractions; the correctness and performance of
particular Lanczos or Spouge implementations; or the zeta functional
equation. The complex-sector form of the full Stirling expansion is
also outside the proved scope: the sharp all-order remainder theorem
here concerns the positive real axis. Binet's formulas do hold on the
certified complex right half-plane as explicitly proved.

Hankel contour representations are not included: proving them would
require an additional finite contour development with branches and
contour-error estimates. Historical statements, specific software
implementations, and long decimal tables are not mathematical theorem
targets of the present argument. The exact certificate run illustrates
the evaluator rather than certifying the source articles' decimal
tables or their quoted continued-fraction digits.

Finally, the first-order claim is the explicit arithmetic theorem schema
for certified primitive recursive presentations. It does not quantify
internally over all classical complex numbers or all real functions.
Nor does the archive claim a completed Lean derivation for all the
listed analysis. These distinctions are part of the statement of the
result, and preserve the constrained constructive setting requested.

\begin{thebibliography}{99}
\bibitem{wiki} Wikipedia contributors, \emph{Gamma function}, accessed 15 November 2026. \url{https://en.wikipedia.org/wiki/Gamma_function}. Used to delimit the requested coverage.
\bibitem{mathworld} E. W. Weisstein, \emph{Gamma Function}, MathWorld, accessed 15 November 2026. \url{https://mathworld.wolfram.com/GammaFunction.html}. Used to delimit the requested coverage.
\bibitem{dlmf} R. A. Askey and R. Roy, \emph{Gamma Function}, Chapter 5 of the NIST Digital Library of Mathematical Functions, especially \S\S5.2, 5.5--5.9, 5.11--5.12, and 5.15. \url{https://dlmf.nist.gov/5}. Formula comparison; the proofs in this report supply their own finite estimates.
\bibitem{repo} V. Reshetnikov, \emph{ProveIt: ConciseIPRA}, repository snapshot \texttt{94a4713736d281e1bfcd0b54e0acb6410136971c}, inspected 15 November 2026. \url{https://github.com/VladimirReshetnikov/ProveIt/tree/94a4713736d281e1bfcd0b54e0acb6410136971c/Logic/PrimitiveRecursiveArithmetic/ConciseIPRA}.
\bibitem{bmrepo} \emph{A Finitary Constructive Bohr--Mollerup Theorem in Intuitionistic Primitive Recursive Arithmetic}, ProveIt draft \texttt{docs/drafts/bohr\_mollerup\_ipra\_2026-11-15.zip}, blob \texttt{55a972a2cd46f9b4a5edd70b79598cca3413ccd7}. Its finite strip brackets and representation discipline inform this development.
\bibitem{stirrepo} \emph{Stirling's formula in intuitionistic primitive recursive arithmetic}, ProveIt draft \texttt{docs/drafts/stirling\_ipra.zip}, blob \texttt{6c6f8e17747d567b426b19f0cb2678cba303457f}. Its finite-difference method is extended here to positive real arguments.
\end{thebibliography}
\end{document}
